%% Generated by Sphinx.
\def\sphinxdocclass{report}
\documentclass[letterpaper,10pt,english]{sphinxmanual}
\ifdefined\pdfpxdimen
   \let\sphinxpxdimen\pdfpxdimen\else\newdimen\sphinxpxdimen
\fi \sphinxpxdimen=.75bp\relax
\ifdefined\pdfimageresolution
    \pdfimageresolution= \numexpr \dimexpr1in\relax/\sphinxpxdimen\relax
\fi
\newdimen\sphinxremdimen\sphinxremdimen = 10pt
%% let collapsible pdf bookmarks panel have high depth per default
\PassOptionsToPackage{bookmarksdepth=5}{hyperref}

\PassOptionsToPackage{booktabs}{sphinx}
\PassOptionsToPackage{colorrows}{sphinx}

\PassOptionsToPackage{warn}{textcomp}
\usepackage[utf8]{inputenc}
\ifdefined\DeclareUnicodeCharacter
% support both utf8 and utf8x syntaxes
  \ifdefined\DeclareUnicodeCharacterAsOptional
    \def\sphinxDUC#1{\DeclareUnicodeCharacter{"#1}}
  \else
    \let\sphinxDUC\DeclareUnicodeCharacter
  \fi
  \sphinxDUC{00A0}{\nobreakspace}
  \sphinxDUC{2500}{\sphinxunichar{2500}}
  \sphinxDUC{2502}{\sphinxunichar{2502}}
  \sphinxDUC{2514}{\sphinxunichar{2514}}
  \sphinxDUC{251C}{\sphinxunichar{251C}}
  \sphinxDUC{2572}{\textbackslash}
\fi
\usepackage{cmap}
\usepackage[T1]{fontenc}
\usepackage{amsmath,amssymb,amstext}
\usepackage{babel}



\usepackage{tgtermes}
\usepackage{tgheros}
\renewcommand{\ttdefault}{txtt}



\usepackage[Bjarne]{fncychap}
\usepackage{sphinx}

\fvset{fontsize=auto}
\usepackage{geometry}


% Include hyperref last.
\usepackage{hyperref}
% Fix anchor placement for figures with captions.
\usepackage{hypcap}% it must be loaded after hyperref.
% Set up styles of URL: it should be placed after hyperref.
\urlstyle{same}


\usepackage{sphinxmessages}
\setcounter{tocdepth}{1}



\title{PsrPopPySuper Documentation}
\date{Jul 06, 2026}
\release{1}
\author{D Tripathi}
\newcommand{\sphinxlogo}{\vbox{}}
\renewcommand{\releasename}{Release}
\makeindex
\begin{document}

\ifdefined\shorthandoff
  \ifnum\catcode`\=\string=\active\shorthandoff{=}\fi
  \ifnum\catcode`\"=\active\shorthandoff{"}\fi
\fi

\pagestyle{empty}
\sphinxmaketitle
\pagestyle{plain}
\sphinxtableofcontents
\pagestyle{normal}
\phantomsection\label{\detokenize{index::doc}}


\sphinxAtStartPar
Contents:

\sphinxstepscope


\chapter{Introduction}
\label{\detokenize{introduction:introduction}}\label{\detokenize{introduction::doc}}
\sphinxAtStartPar
PsrPopPySuper is a modernized Python3 port of Sam Bates’ PsrPopPy code (\sphinxurl{https://github.com/samb8s/psrpoppy}).
This project retains the original code’s functionality and most of the user facing APIs while updating interfaces and adding newer
electron models like NE2025 and YMW16.

\sphinxAtStartPar
Contributions are welcome via the project’s GitHub repository.

\sphinxstepscope


\chapter{Installation}
\label{\detokenize{installation:installation}}\label{\detokenize{installation:id1}}\label{\detokenize{installation::doc}}
\sphinxAtStartPar
To get started with PsrPopPySuper there are a few steps you’ll need to
go through.

\sphinxAtStartPar
PsrPopPySuper is currently supported on Linux and Mac OS X, with python
version 3.11 or later.


\section{Download the package}
\label{\detokenize{installation:download-the-package}}\label{\detokenize{installation:download-package}}
\sphinxAtStartPar
The source for PsrPopPy can be downloaded from \sphinxhref{http://github.com}{GitHub}
from the \sphinxhref{https://github.com/samb8s/PsrPopPy}{PsrPopPy page}.
The source will contain the Python modules and scripts needed both for
basic and advanced use.


\section{Installing the package}
\label{\detokenize{installation:installing-the-package}}\label{\detokenize{installation:installing-package}}
\sphinxAtStartPar
PsrPopPy now supports standard Python package installation using pip.
This automatically handles the Fortran compilation and installation.


\subsection{Prerequisites}
\label{\detokenize{installation:prerequisites}}\begin{itemize}
\item {} 
\sphinxAtStartPar
Python 3.11 or later

\item {} 
\sphinxAtStartPar
gfortran compiler (for Fortran library compilation)

\end{itemize}

\sphinxAtStartPar
On macOS, install gfortran using Homebrew:

\begin{sphinxVerbatim}[commandchars=\\\{\}]
\PYG{n}{brew} \PYG{n}{install} \PYG{n}{gcc}
\end{sphinxVerbatim}

\sphinxAtStartPar
On Linux, install gfortran using your package manager:

\begin{sphinxVerbatim}[commandchars=\\\{\}]
\PYG{c+c1}{\PYGZsh{} Ubuntu/Debian}
\PYG{n}{sudo} \PYG{n}{apt}\PYG{o}{\PYGZhy{}}\PYG{n}{get} \PYG{n}{install} \PYG{n}{gfortran}

\PYG{c+c1}{\PYGZsh{} CentOS/RHEL/Fedora}
\PYG{n}{sudo} \PYG{n}{yum} \PYG{n}{install} \PYG{n}{gcc}\PYG{o}{\PYGZhy{}}\PYG{n}{gfortran}  \PYG{c+c1}{\PYGZsh{} or dnf install gcc\PYGZhy{}gfortran}
\end{sphinxVerbatim}


\subsection{Installation}
\label{\detokenize{installation:id2}}

\subsubsection{With Conda}
\label{\detokenize{installation:with-conda}}
\sphinxAtStartPar
Use the {[}environment.yml{]}(environment.yml) to make the conda environment and activate it:

\begin{sphinxVerbatim}[commandchars=\\\{\}]
\PYG{n}{conda} \PYG{n}{create} \PYG{o}{\PYGZhy{}}\PYG{o}{\PYGZhy{}}\PYG{n}{file} \PYG{n}{environment}\PYG{o}{.}\PYG{n}{yml} \PYG{o}{\PYGZhy{}}\PYG{o}{\PYGZhy{}}\PYG{n}{name} \PYG{n}{PsrPopPySuper}
\PYG{n}{conda} \PYG{n}{activate} \PYG{n}{PsrpoppySuper}
\end{sphinxVerbatim}

\sphinxAtStartPar
Replace PsrPopPySuper with your own preferred name of the environment:

\begin{sphinxVerbatim}[commandchars=\\\{\}]
\PYG{n}{pip} \PYG{n}{install} \PYG{o}{\PYGZhy{}}\PYG{n}{e} \PYG{o}{.}
\end{sphinxVerbatim}


\subsubsection{Without Conda}
\label{\detokenize{installation:without-conda}}
\sphinxAtStartPar
Before moving on without conda ensure you have python version 3.11+ and have working c and fortran compiler:

\begin{sphinxVerbatim}[commandchars=\\\{\}]
\PYG{n}{cd} \PYG{n}{PsrPopPySuper}
\PYG{n}{pip} \PYG{n}{install} \PYG{o}{\PYGZhy{}}\PYG{n}{r} \PYG{n}{requirements}\PYG{o}{.}\PYG{n}{txt}
\PYG{n}{pip} \PYG{n}{install} \PYG{o}{\PYGZhy{}}\PYG{n}{e} \PYG{o}{.}
\end{sphinxVerbatim}

\sphinxAtStartPar
{\color{red}\bfseries{}\textasciigrave{}}

\sphinxAtStartPar
This will:
\begin{itemize}
\item {} 
\sphinxAtStartPar
Compile the Fortran and C libraries using gfortran and gcc

\item {} 
\sphinxAtStartPar
Install the Python package

\item {} 
\sphinxAtStartPar
Install command\sphinxhyphen{}line scripts (dosurvey, populate, evolve)

\end{itemize}


\section{Legacy Fortran Compilation (if needed)}
\label{\detokenize{installation:legacy-fortran-compilation-if-needed}}\label{\detokenize{installation:legacy-compiling-fortran}}
\sphinxAtStartPar
Although PsrPopPySuper is a Python\sphinxhyphen{}based package, some of the algorithms
have been kept in their native FORTRAN for speed and ease of
programming (e.g. the NE2001 electron model, coordinate conversion…).
The pip installation now handles this automatically, but if you need
to compile manually:

\sphinxAtStartPar
From the base directory:

\begin{sphinxVerbatim}[commandchars=\\\{\}]
\PYG{n}{cd} \PYG{n}{psrpoppy}\PYG{o}{/}\PYG{n}{fortran}
\end{sphinxVerbatim}

\sphinxAtStartPar
To use \sphinxcode{\sphinxupquote{make}}, edit \sphinxcode{\sphinxupquote{makefile.\textless{}OSTYPE\textgreater{}}} and ensure that the gf variable
points to the location of your gfortran compiler. Then simply type \sphinxcode{\sphinxupquote{make}}.
All being well, four .so files will be generated.

\sphinxAtStartPar
Failing this, edit either \sphinxcode{\sphinxupquote{make\_mac.sh}} or \sphinxcode{\sphinxupquote{make\_linux.sh}}, depending upon
your system, so that the \sphinxcode{\sphinxupquote{gf}} variable points to your local gfortran/f77
compiler. Running the script:

\begin{sphinxVerbatim}[commandchars=\\\{\}]
\PYG{n}{bash} \PYG{n}{make\PYGZus{}}\PYG{o}{\PYGZlt{}}\PYG{n}{os}\PYG{o}{\PYGZgt{}}\PYG{o}{.}\PYG{n}{sh}
\end{sphinxVerbatim}

\sphinxstepscope


\chapter{Command\sphinxhyphen{}line scripts}
\label{\detokenize{cmd-line-docs:command-line-scripts}}\label{\detokenize{cmd-line-docs:command-line-docs}}\label{\detokenize{cmd-line-docs::doc}}

\section{populate.py}
\label{\detokenize{cmd-line-docs:populate-py}}\label{\detokenize{cmd-line-docs:populate-docs}}\index{populate.py command line option@\spxentry{populate.py command line option}!\sphinxhyphen{}n@\spxentry{\sphinxhyphen{}n}}\index{\sphinxhyphen{}n@\spxentry{\sphinxhyphen{}n}!populate.py command line option@\spxentry{populate.py command line option}}

\begin{fulllineitems}
\phantomsection\label{\detokenize{cmd-line-docs:cmdoption-populate.py-n}}
\pysigstartsignatures
\pysigline
{\sphinxbfcode{\sphinxupquote{\sphinxhyphen{}n}}\sphinxcode{\sphinxupquote{~\textless{}number~of~pulsars\textgreater{}}}}
\pysigstopsignatures
\sphinxAtStartPar
Required: Number of pulsars to generate; or to detect in a survey

\end{fulllineitems}

\index{populate.py command line option@\spxentry{populate.py command line option}!\sphinxhyphen{}o@\spxentry{\sphinxhyphen{}o}}\index{\sphinxhyphen{}o@\spxentry{\sphinxhyphen{}o}!populate.py command line option@\spxentry{populate.py command line option}}

\begin{fulllineitems}
\phantomsection\label{\detokenize{cmd-line-docs:cmdoption-populate.py-o}}
\pysigstartsignatures
\pysigline
{\sphinxbfcode{\sphinxupquote{\sphinxhyphen{}o}}\sphinxcode{\sphinxupquote{~\textless{}output\textgreater{}}}}
\pysigstopsignatures
\sphinxAtStartPar
Output file name for population model (def=populate.model)

\end{fulllineitems}

\index{populate.py command line option@\spxentry{populate.py command line option}!\sphinxhyphen{}asc@\spxentry{\sphinxhyphen{}asc}}\index{\sphinxhyphen{}asc@\spxentry{\sphinxhyphen{}asc}!populate.py command line option@\spxentry{populate.py command line option}}

\begin{fulllineitems}
\phantomsection\label{\detokenize{cmd-line-docs:cmdoption-populate.py-asc}}
\pysigstartsignatures
\pysigline
{\sphinxbfcode{\sphinxupquote{\sphinxhyphen{}asc}}\sphinxcode{\sphinxupquote{~\textless{}ascii~output~file~name\textgreater{}}}}
\pysigstopsignatures
\sphinxAtStartPar
Output file name for ascii output file (def=populate.txt)

\end{fulllineitems}

\index{populate.py command line option@\spxentry{populate.py command line option}!\sphinxhyphen{}surveys@\spxentry{\sphinxhyphen{}surveys}}\index{\sphinxhyphen{}surveys@\spxentry{\sphinxhyphen{}surveys}!populate.py command line option@\spxentry{populate.py command line option}}

\begin{fulllineitems}
\phantomsection\label{\detokenize{cmd-line-docs:cmdoption-populate.py-surveys}}
\pysigstartsignatures
\pysigline
{\sphinxbfcode{\sphinxupquote{\sphinxhyphen{}surveys}}\sphinxcode{\sphinxupquote{~\textless{}SURVEY~NAME(S)\textgreater{}}}}
\pysigstopsignatures
\sphinxAtStartPar
List of surveys to use when trying to detect pulsars (default=None)

\end{fulllineitems}

\index{populate.py command line option@\spxentry{populate.py command line option}!\sphinxhyphen{}z@\spxentry{\sphinxhyphen{}z}}\index{\sphinxhyphen{}z@\spxentry{\sphinxhyphen{}z}!populate.py command line option@\spxentry{populate.py command line option}}

\begin{fulllineitems}
\phantomsection\label{\detokenize{cmd-line-docs:cmdoption-populate.py-z}}
\pysigstartsignatures
\pysigline
{\sphinxbfcode{\sphinxupquote{\sphinxhyphen{}z}}\sphinxcode{\sphinxupquote{~\textless{}scale~height\textgreater{}}}}
\pysigstopsignatures
\sphinxAtStartPar
Scale height of pulsars about Galactic plane, in kpc (default=0.33)

\end{fulllineitems}

\index{populate.py command line option@\spxentry{populate.py command line option}!\sphinxhyphen{}w@\spxentry{\sphinxhyphen{}w}}\index{\sphinxhyphen{}w@\spxentry{\sphinxhyphen{}w}!populate.py command line option@\spxentry{populate.py command line option}}

\begin{fulllineitems}
\phantomsection\label{\detokenize{cmd-line-docs:cmdoption-populate.py-w}}
\pysigstartsignatures
\pysigline
{\sphinxbfcode{\sphinxupquote{\sphinxhyphen{}w}}\sphinxcode{\sphinxupquote{~\textless{}width\textgreater{}}}}
\pysigstopsignatures
\sphinxAtStartPar
Pulse width to use when generating pulsars (default=0, use beaming model)

\end{fulllineitems}

\index{populate.py command line option@\spxentry{populate.py command line option}!\sphinxhyphen{}si@\spxentry{\sphinxhyphen{}si}}\index{\sphinxhyphen{}si@\spxentry{\sphinxhyphen{}si}!populate.py command line option@\spxentry{populate.py command line option}}

\begin{fulllineitems}
\phantomsection\label{\detokenize{cmd-line-docs:cmdoption-populate.py-si}}
\pysigstartsignatures
\pysigline
{\sphinxbfcode{\sphinxupquote{\sphinxhyphen{}si}}\sphinxcode{\sphinxupquote{~\textless{}SImean~SIsigma\textgreater{}}}}
\pysigstopsignatures
\sphinxAtStartPar
Spectral index mean and standard deviation (default=\sphinxhyphen{}1.6, 0.35)

\end{fulllineitems}

\index{populate.py command line option@\spxentry{populate.py command line option}!\sphinxhyphen{}sc@\spxentry{\sphinxhyphen{}sc}}\index{\sphinxhyphen{}sc@\spxentry{\sphinxhyphen{}sc}!populate.py command line option@\spxentry{populate.py command line option}}

\begin{fulllineitems}
\phantomsection\label{\detokenize{cmd-line-docs:cmdoption-populate.py-sc}}
\pysigstartsignatures
\pysigline
{\sphinxbfcode{\sphinxupquote{\sphinxhyphen{}sc}}\sphinxcode{\sphinxupquote{~\textless{}scatter~index\textgreater{}}}}
\pysigstopsignatures
\sphinxAtStartPar
Spectral index of scattering law to use (default=\sphinxhyphen{}3.86, Bhat et al model)

\end{fulllineitems}

\index{populate.py command line option@\spxentry{populate.py command line option}!\sphinxhyphen{}pdist@\spxentry{\sphinxhyphen{}pdist}}\index{\sphinxhyphen{}pdist@\spxentry{\sphinxhyphen{}pdist}!populate.py command line option@\spxentry{populate.py command line option}}

\begin{fulllineitems}
\phantomsection\label{\detokenize{cmd-line-docs:cmdoption-populate.py-pdist}}
\pysigstartsignatures
\pysigline
{\sphinxbfcode{\sphinxupquote{\sphinxhyphen{}pdist}}\sphinxcode{\sphinxupquote{~\textless{}distribution~name\textgreater{}}}}
\pysigstopsignatures
\sphinxAtStartPar
Distribution type for pulse periods (default=lnorm)

\sphinxAtStartPar
Supported: ‘lnorm’, ‘norm’, ‘cc97’

\end{fulllineitems}

\index{populate.py command line option@\spxentry{populate.py command line option}!\sphinxhyphen{}p@\spxentry{\sphinxhyphen{}p}}\index{\sphinxhyphen{}p@\spxentry{\sphinxhyphen{}p}!populate.py command line option@\spxentry{populate.py command line option}}

\begin{fulllineitems}
\phantomsection\label{\detokenize{cmd-line-docs:cmdoption-populate.py-p}}
\pysigstartsignatures
\pysigline
{\sphinxbfcode{\sphinxupquote{\sphinxhyphen{}p}}\sphinxcode{\sphinxupquote{~\textless{}mean~stddev\textgreater{}}}}
\pysigstopsignatures
\sphinxAtStartPar
Mean and standard deviation to use in period dist ‘lnorm’ or ‘norm’ (def=\sphinxhyphen{}2.7, 0.34)

\end{fulllineitems}

\index{populate.py command line option@\spxentry{populate.py command line option}!\sphinxhyphen{}rdist@\spxentry{\sphinxhyphen{}rdist}}\index{\sphinxhyphen{}rdist@\spxentry{\sphinxhyphen{}rdist}!populate.py command line option@\spxentry{populate.py command line option}}

\begin{fulllineitems}
\phantomsection\label{\detokenize{cmd-line-docs:cmdoption-populate.py-rdist}}
\pysigstartsignatures
\pysigline
{\sphinxbfcode{\sphinxupquote{\sphinxhyphen{}rdist}}\sphinxcode{\sphinxupquote{~\textless{}radial~model\textgreater{}}}}
\pysigstopsignatures
\sphinxAtStartPar
Model to use for Galactic radial distribution of pulsars

\sphinxAtStartPar
Supported: ‘lfl06’, ‘yk04’, ‘isotropic’, ‘slab’, ‘disk’

\end{fulllineitems}

\index{populate.py command line option@\spxentry{populate.py command line option}!\sphinxhyphen{}dm@\spxentry{\sphinxhyphen{}dm}}\index{\sphinxhyphen{}dm@\spxentry{\sphinxhyphen{}dm}!populate.py command line option@\spxentry{populate.py command line option}}

\begin{fulllineitems}
\phantomsection\label{\detokenize{cmd-line-docs:cmdoption-populate.py-dm}}
\pysigstartsignatures
\pysigline
{\sphinxbfcode{\sphinxupquote{\sphinxhyphen{}dm}}\sphinxcode{\sphinxupquote{~\textless{}Electron~model\textgreater{}}}}
\pysigstopsignatures
\sphinxAtStartPar
Model to describe the Galactic electron distribution

\sphinxAtStartPar
Supported: ‘ne2025’, ‘ne2001’, ‘ymw16’

\end{fulllineitems}

\index{populate.py command line option@\spxentry{populate.py command line option}!\sphinxhyphen{}gps@\spxentry{\sphinxhyphen{}gps}}\index{\sphinxhyphen{}gps@\spxentry{\sphinxhyphen{}gps}!populate.py command line option@\spxentry{populate.py command line option}}

\begin{fulllineitems}
\phantomsection\label{\detokenize{cmd-line-docs:cmdoption-populate.py-gps}}
\pysigstartsignatures
\pysigline
{\sphinxbfcode{\sphinxupquote{\sphinxhyphen{}gps}}\sphinxcode{\sphinxupquote{~\textless{}fraction~\textquotesingle{}a\textquotesingle{}\textgreater{}}}}
\pysigstopsignatures
\sphinxAtStartPar
Add \textless{}fraction\textgreater{} pulsars with GHz\sphinxhyphen{}frequency turnovers with index a

\end{fulllineitems}

\index{populate.py command line option@\spxentry{populate.py command line option}!\sphinxhyphen{}doublespec@\spxentry{\sphinxhyphen{}doublespec}}\index{\sphinxhyphen{}doublespec@\spxentry{\sphinxhyphen{}doublespec}!populate.py command line option@\spxentry{populate.py command line option}}

\begin{fulllineitems}
\phantomsection\label{\detokenize{cmd-line-docs:cmdoption-populate.py-doublespec}}
\pysigstartsignatures
\pysigline
{\sphinxbfcode{\sphinxupquote{\sphinxhyphen{}doublespec}}\sphinxcode{\sphinxupquote{~\textless{}fraction~alpha1\textgreater{}}}}
\pysigstopsignatures
\sphinxAtStartPar
Add \textless{}fraction\textgreater{} pulsars with low\sphinxhyphen{}frequency (below 1GHz) spectral index of alpha1

\end{fulllineitems}

\index{populate.py command line option@\spxentry{populate.py command line option}!\sphinxhyphen{}\sphinxhyphen{}nostdout@\spxentry{\sphinxhyphen{}\sphinxhyphen{}nostdout}}\index{\sphinxhyphen{}\sphinxhyphen{}nostdout@\spxentry{\sphinxhyphen{}\sphinxhyphen{}nostdout}!populate.py command line option@\spxentry{populate.py command line option}}

\begin{fulllineitems}
\phantomsection\label{\detokenize{cmd-line-docs:cmdoption-populate.py-nostdout}}
\pysigstartsignatures
\pysigline
{\sphinxbfcode{\sphinxupquote{\sphinxhyphen{}\sphinxhyphen{}nostdout}}\sphinxcode{\sphinxupquote{}}}
\pysigstopsignatures
\sphinxAtStartPar
Turn off writing to stdout. Useful for many iterations eg. in large simulations

\end{fulllineitems}



\section{dosurvey.py}
\label{\detokenize{cmd-line-docs:dosurvey-py}}\label{\detokenize{cmd-line-docs:dosurvey-docs}}\index{dosurvey.py command line option@\spxentry{dosurvey.py command line option}!\sphinxhyphen{}f@\spxentry{\sphinxhyphen{}f}}\index{\sphinxhyphen{}f@\spxentry{\sphinxhyphen{}f}!dosurvey.py command line option@\spxentry{dosurvey.py command line option}}

\begin{fulllineitems}
\phantomsection\label{\detokenize{cmd-line-docs:cmdoption-dosurvey.py-f}}
\pysigstartsignatures
\pysigline
{\sphinxbfcode{\sphinxupquote{\sphinxhyphen{}f}}\sphinxcode{\sphinxupquote{~\textless{}filename\textgreater{}}}}
\pysigstopsignatures
\sphinxAtStartPar
Input population model to use (default=populate.model)

\end{fulllineitems}

\index{dosurvey.py command line option@\spxentry{dosurvey.py command line option}!\sphinxhyphen{}surveys@\spxentry{\sphinxhyphen{}surveys}}\index{\sphinxhyphen{}surveys@\spxentry{\sphinxhyphen{}surveys}!dosurvey.py command line option@\spxentry{dosurvey.py command line option}}

\begin{fulllineitems}
\phantomsection\label{\detokenize{cmd-line-docs:cmdoption-dosurvey.py-surveys}}
\pysigstartsignatures
\pysigline
{\sphinxbfcode{\sphinxupquote{\sphinxhyphen{}surveys}}\sphinxcode{\sphinxupquote{~\textless{}SURVEY~NAME(S)\textgreater{}}}}
\pysigstopsignatures
\sphinxAtStartPar
Required: Name(s) of surveys to simulate on the population model

\end{fulllineitems}

\index{dosurvey.py command line option@\spxentry{dosurvey.py command line option}!\sphinxhyphen{}\sphinxhyphen{}noresults@\spxentry{\sphinxhyphen{}\sphinxhyphen{}noresults}}\index{\sphinxhyphen{}\sphinxhyphen{}noresults@\spxentry{\sphinxhyphen{}\sphinxhyphen{}noresults}!dosurvey.py command line option@\spxentry{dosurvey.py command line option}}

\begin{fulllineitems}
\phantomsection\label{\detokenize{cmd-line-docs:cmdoption-dosurvey.py-noresults}}
\pysigstartsignatures
\pysigline
{\sphinxbfcode{\sphinxupquote{\sphinxhyphen{}\sphinxhyphen{}noresults}}\sphinxcode{\sphinxupquote{}}}
\pysigstopsignatures
\sphinxAtStartPar
By default, a .results file is written, containing a model of the population
detected in the survey. This option switches off the writing of this file.

\end{fulllineitems}

\index{dosurvey.py command line option@\spxentry{dosurvey.py command line option}!\sphinxhyphen{}\sphinxhyphen{}asc@\spxentry{\sphinxhyphen{}\sphinxhyphen{}asc}}\index{\sphinxhyphen{}\sphinxhyphen{}asc@\spxentry{\sphinxhyphen{}\sphinxhyphen{}asc}!dosurvey.py command line option@\spxentry{dosurvey.py command line option}}

\begin{fulllineitems}
\phantomsection\label{\detokenize{cmd-line-docs:cmdoption-dosurvey.py-asc}}
\pysigstartsignatures
\pysigline
{\sphinxbfcode{\sphinxupquote{\sphinxhyphen{}\sphinxhyphen{}asc}}\sphinxcode{\sphinxupquote{}}}
\pysigstopsignatures
\sphinxAtStartPar
Write the survey model in plain ascii (psrpop old style). Not recommended,
since the cPickle ‘.results’ file is easier to work with.

\end{fulllineitems}

\index{dosurvey.py command line option@\spxentry{dosurvey.py command line option}!\sphinxhyphen{}\sphinxhyphen{}summary@\spxentry{\sphinxhyphen{}\sphinxhyphen{}summary}}\index{\sphinxhyphen{}\sphinxhyphen{}summary@\spxentry{\sphinxhyphen{}\sphinxhyphen{}summary}!dosurvey.py command line option@\spxentry{dosurvey.py command line option}}

\begin{fulllineitems}
\phantomsection\label{\detokenize{cmd-line-docs:cmdoption-dosurvey.py-summary}}
\pysigstartsignatures
\pysigline
{\sphinxbfcode{\sphinxupquote{\sphinxhyphen{}\sphinxhyphen{}summary}}\sphinxcode{\sphinxupquote{}}}
\pysigstopsignatures
\sphinxAtStartPar
Write a short .summary file (per survey) describing number of detections,
number of pulsars outside survey area, number smeared, and number not beaming

\end{fulllineitems}

\index{dosurvey.py command line option@\spxentry{dosurvey.py command line option}!\sphinxhyphen{}\sphinxhyphen{}nostdout@\spxentry{\sphinxhyphen{}\sphinxhyphen{}nostdout}}\index{\sphinxhyphen{}\sphinxhyphen{}nostdout@\spxentry{\sphinxhyphen{}\sphinxhyphen{}nostdout}!dosurvey.py command line option@\spxentry{dosurvey.py command line option}}

\begin{fulllineitems}
\phantomsection\label{\detokenize{cmd-line-docs:cmdoption-dosurvey.py-nostdout}}
\pysigstartsignatures
\pysigline
{\sphinxbfcode{\sphinxupquote{\sphinxhyphen{}\sphinxhyphen{}nostdout}}\sphinxcode{\sphinxupquote{}}}
\pysigstopsignatures
\sphinxAtStartPar
Turn off writing to stdout. Useful for many iterations eg. in large simulations

\end{fulllineitems}



\section{evolve.py}
\label{\detokenize{cmd-line-docs:evolve-py}}\label{\detokenize{cmd-line-docs:evolve-docs}}\index{evolve.py command line option@\spxentry{evolve.py command line option}!\sphinxhyphen{}n@\spxentry{\sphinxhyphen{}n}}\index{\sphinxhyphen{}n@\spxentry{\sphinxhyphen{}n}!evolve.py command line option@\spxentry{evolve.py command line option}}

\begin{fulllineitems}
\phantomsection\label{\detokenize{cmd-line-docs:cmdoption-evolve.py-n}}
\pysigstartsignatures
\pysigline
{\sphinxbfcode{\sphinxupquote{\sphinxhyphen{}n}}\sphinxcode{\sphinxupquote{~\textless{}number~of~pulsars\textgreater{}}}}
\pysigstopsignatures
\sphinxAtStartPar
Required: Number of pulsars to generate or detect

\end{fulllineitems}

\index{evolve.py command line option@\spxentry{evolve.py command line option}!\sphinxhyphen{}o@\spxentry{\sphinxhyphen{}o}}\index{\sphinxhyphen{}o@\spxentry{\sphinxhyphen{}o}!evolve.py command line option@\spxentry{evolve.py command line option}}

\begin{fulllineitems}
\phantomsection\label{\detokenize{cmd-line-docs:cmdoption-evolve.py-o}}
\pysigstartsignatures
\pysigline
{\sphinxbfcode{\sphinxupquote{\sphinxhyphen{}o}}\sphinxcode{\sphinxupquote{~\textless{}output\textgreater{}}}}
\pysigstopsignatures
\sphinxAtStartPar
Output file name for population model (def=evolve.model)

\end{fulllineitems}

\index{evolve.py command line option@\spxentry{evolve.py command line option}!\sphinxhyphen{}asc@\spxentry{\sphinxhyphen{}asc}}\index{\sphinxhyphen{}asc@\spxentry{\sphinxhyphen{}asc}!evolve.py command line option@\spxentry{evolve.py command line option}}

\begin{fulllineitems}
\phantomsection\label{\detokenize{cmd-line-docs:cmdoption-evolve.py-asc}}
\pysigstartsignatures
\pysigline
{\sphinxbfcode{\sphinxupquote{\sphinxhyphen{}asc}}\sphinxcode{\sphinxupquote{~\textless{}ascii~output~file~name\textgreater{}}}}
\pysigstopsignatures
\sphinxAtStartPar
Output file name for ascii output file (def=None)

\end{fulllineitems}

\index{evolve.py command line option@\spxentry{evolve.py command line option}!\sphinxhyphen{}surveys@\spxentry{\sphinxhyphen{}surveys}}\index{\sphinxhyphen{}surveys@\spxentry{\sphinxhyphen{}surveys}!evolve.py command line option@\spxentry{evolve.py command line option}}

\begin{fulllineitems}
\phantomsection\label{\detokenize{cmd-line-docs:cmdoption-evolve.py-surveys}}
\pysigstartsignatures
\pysigline
{\sphinxbfcode{\sphinxupquote{\sphinxhyphen{}surveys}}\sphinxcode{\sphinxupquote{~\textless{}SURVEY~NAME(S)\textgreater{}}}}
\pysigstopsignatures
\sphinxAtStartPar
List of surveys to use when trying to detect pulsars (default=None)

\end{fulllineitems}

\index{evolve.py command line option@\spxentry{evolve.py command line option}!\sphinxhyphen{}dm@\spxentry{\sphinxhyphen{}dm}}\index{\sphinxhyphen{}dm@\spxentry{\sphinxhyphen{}dm}!evolve.py command line option@\spxentry{evolve.py command line option}}

\begin{fulllineitems}
\phantomsection\label{\detokenize{cmd-line-docs:cmdoption-evolve.py-dm}}
\pysigstartsignatures
\pysigline
{\sphinxbfcode{\sphinxupquote{\sphinxhyphen{}dm}}\sphinxcode{\sphinxupquote{~\textless{}Electron~model\textgreater{}}}}
\pysigstopsignatures
\sphinxAtStartPar
Model to describe the Galactic electron distribution

\sphinxAtStartPar
Supported: ‘ne2025’, , ‘ne2001’, ‘ymw16’

\end{fulllineitems}

\index{evolve.py command line option@\spxentry{evolve.py command line option}!\sphinxhyphen{}tmax@\spxentry{\sphinxhyphen{}tmax}}\index{\sphinxhyphen{}tmax@\spxentry{\sphinxhyphen{}tmax}!evolve.py command line option@\spxentry{evolve.py command line option}}

\begin{fulllineitems}
\phantomsection\label{\detokenize{cmd-line-docs:cmdoption-evolve.py-tmax}}
\pysigstartsignatures
\pysigline
{\sphinxbfcode{\sphinxupquote{\sphinxhyphen{}tmax}}\sphinxcode{\sphinxupquote{~\textless{}max~age\textgreater{}}}}
\pysigstopsignatures
\sphinxAtStartPar
Maximum initial age of pulsars, in years (def=1.0E9)

\end{fulllineitems}

\index{evolve.py command line option@\spxentry{evolve.py command line option}!\sphinxhyphen{}p@\spxentry{\sphinxhyphen{}p}}\index{\sphinxhyphen{}p@\spxentry{\sphinxhyphen{}p}!evolve.py command line option@\spxentry{evolve.py command line option}}

\begin{fulllineitems}
\phantomsection\label{\detokenize{cmd-line-docs:cmdoption-evolve.py-p}}
\pysigstartsignatures
\pysigline
{\sphinxbfcode{\sphinxupquote{\sphinxhyphen{}p}}\sphinxcode{\sphinxupquote{~\textless{}mean~stddev\textgreater{}}}}
\pysigstopsignatures
\sphinxAtStartPar
Period distribution mean and standard deviation, seconds (def=0.3, 0.15)

\end{fulllineitems}

\index{evolve.py command line option@\spxentry{evolve.py command line option}!\sphinxhyphen{}ldist@\spxentry{\sphinxhyphen{}ldist}}\index{\sphinxhyphen{}ldist@\spxentry{\sphinxhyphen{}ldist}!evolve.py command line option@\spxentry{evolve.py command line option}}

\begin{fulllineitems}
\phantomsection\label{\detokenize{cmd-line-docs:cmdoption-evolve.py-ldist}}
\pysigstartsignatures
\pysigline
{\sphinxbfcode{\sphinxupquote{\sphinxhyphen{}ldist}}\sphinxcode{\sphinxupquote{~\textless{}distribution\textgreater{}}}}
\pysigstopsignatures
\sphinxAtStartPar
Distribution to use for luminosities

\sphinxAtStartPar
Supported: ‘fk06’, ‘lnorm’, ‘pow’

\end{fulllineitems}

\index{evolve.py command line option@\spxentry{evolve.py command line option}!\sphinxhyphen{}l@\spxentry{\sphinxhyphen{}l}}\index{\sphinxhyphen{}l@\spxentry{\sphinxhyphen{}l}!evolve.py command line option@\spxentry{evolve.py command line option}}

\begin{fulllineitems}
\phantomsection\label{\detokenize{cmd-line-docs:cmdoption-evolve.py-l}}
\pysigstartsignatures
\pysigline
{\sphinxbfcode{\sphinxupquote{\sphinxhyphen{}l}}\sphinxcode{\sphinxupquote{~\textless{}parameters\textgreater{}}}}
\pysigstopsignatures
\sphinxAtStartPar
Luminosity distribution parameters (def=\sphinxhyphen{}1.5, 0.5)

\end{fulllineitems}

\index{evolve.py command line option@\spxentry{evolve.py command line option}!\sphinxhyphen{}b@\spxentry{\sphinxhyphen{}b}}\index{\sphinxhyphen{}b@\spxentry{\sphinxhyphen{}b}!evolve.py command line option@\spxentry{evolve.py command line option}}

\begin{fulllineitems}
\phantomsection\label{\detokenize{cmd-line-docs:cmdoption-evolve.py-b}}
\pysigstartsignatures
\pysigline
{\sphinxbfcode{\sphinxupquote{\sphinxhyphen{}b}}\sphinxcode{\sphinxupquote{~\textless{}mean~stddev\textgreater{}}}}
\pysigstopsignatures
\sphinxAtStartPar
Mean and standard deviation of log\sphinxhyphen{}normal B field distribution (Gauss, def=12.65, 0.55)

\end{fulllineitems}

\index{evolve.py command line option@\spxentry{evolve.py command line option}!\sphinxhyphen{}vmodel@\spxentry{\sphinxhyphen{}vmodel}}\index{\sphinxhyphen{}vmodel@\spxentry{\sphinxhyphen{}vmodel}!evolve.py command line option@\spxentry{evolve.py command line option}}

\begin{fulllineitems}
\phantomsection\label{\detokenize{cmd-line-docs:cmdoption-evolve.py-vmodel}}
\pysigstartsignatures
\pysigline
{\sphinxbfcode{\sphinxupquote{\sphinxhyphen{}vmodel}}\sphinxcode{\sphinxupquote{~\textless{}model\textgreater{}}}}
\pysigstopsignatures
\sphinxAtStartPar
Velocity model to use

\sphinxAtStartPar
Supported: ‘gaussian’, ‘exp’

\end{fulllineitems}

\index{evolve.py command line option@\spxentry{evolve.py command line option}!\sphinxhyphen{}v@\spxentry{\sphinxhyphen{}v}}\index{\sphinxhyphen{}v@\spxentry{\sphinxhyphen{}v}!evolve.py command line option@\spxentry{evolve.py command line option}}

\begin{fulllineitems}
\phantomsection\label{\detokenize{cmd-line-docs:cmdoption-evolve.py-v}}
\pysigstartsignatures
\pysigline
{\sphinxbfcode{\sphinxupquote{\sphinxhyphen{}v}}\sphinxcode{\sphinxupquote{~\textless{}mean~stddev\textgreater{}}}}
\pysigstopsignatures
\sphinxAtStartPar
Velocity distribution values (def=0, 180 km/s)

\end{fulllineitems}

\index{evolve.py command line option@\spxentry{evolve.py command line option}!\sphinxhyphen{}si@\spxentry{\sphinxhyphen{}si}}\index{\sphinxhyphen{}si@\spxentry{\sphinxhyphen{}si}!evolve.py command line option@\spxentry{evolve.py command line option}}

\begin{fulllineitems}
\phantomsection\label{\detokenize{cmd-line-docs:cmdoption-evolve.py-si}}
\pysigstartsignatures
\pysigline
{\sphinxbfcode{\sphinxupquote{\sphinxhyphen{}si}}\sphinxcode{\sphinxupquote{~\textless{}mean~stddev\textgreater{}}}}
\pysigstopsignatures
\sphinxAtStartPar
Spectral index mean and standard deviation (def=\sphinxhyphen{}1.4, 0.96)

\end{fulllineitems}

\index{evolve.py command line option@\spxentry{evolve.py command line option}!\sphinxhyphen{}spinmodel@\spxentry{\sphinxhyphen{}spinmodel}}\index{\sphinxhyphen{}spinmodel@\spxentry{\sphinxhyphen{}spinmodel}!evolve.py command line option@\spxentry{evolve.py command line option}}

\begin{fulllineitems}
\phantomsection\label{\detokenize{cmd-line-docs:cmdoption-evolve.py-spinmodel}}
\pysigstartsignatures
\pysigline
{\sphinxbfcode{\sphinxupquote{\sphinxhyphen{}spinmodel}}\sphinxcode{\sphinxupquote{~\textless{}model\textgreater{}}}}
\pysigstopsignatures
\sphinxAtStartPar
Spin\sphinxhyphen{}down model to employ

\sphinxAtStartPar
Supported: ‘fk06’, ‘cs06’

\end{fulllineitems}

\index{evolve.py command line option@\spxentry{evolve.py command line option}!\sphinxhyphen{}alignmodel@\spxentry{\sphinxhyphen{}alignmodel}}\index{\sphinxhyphen{}alignmodel@\spxentry{\sphinxhyphen{}alignmodel}!evolve.py command line option@\spxentry{evolve.py command line option}}

\begin{fulllineitems}
\phantomsection\label{\detokenize{cmd-line-docs:cmdoption-evolve.py-alignmodel}}
\pysigstartsignatures
\pysigline
{\sphinxbfcode{\sphinxupquote{\sphinxhyphen{}alignmodel}}\sphinxcode{\sphinxupquote{~\textless{}model\textgreater{}}}}
\pysigstopsignatures
\sphinxAtStartPar
Pulsar alignment model to use

\sphinxAtStartPar
Supported: ‘orthogonal’, ‘random’, ‘rand45’, ‘wj08’

\end{fulllineitems}

\index{evolve.py command line option@\spxentry{evolve.py command line option}!\sphinxhyphen{}aligntime@\spxentry{\sphinxhyphen{}aligntime}}\index{\sphinxhyphen{}aligntime@\spxentry{\sphinxhyphen{}aligntime}!evolve.py command line option@\spxentry{evolve.py command line option}}

\begin{fulllineitems}
\phantomsection\label{\detokenize{cmd-line-docs:cmdoption-evolve.py-aligntime}}
\pysigstartsignatures
\pysigline
{\sphinxbfcode{\sphinxupquote{\sphinxhyphen{}aligntime}}\sphinxcode{\sphinxupquote{~\textless{}time\textgreater{}}}}
\pysigstopsignatures
\sphinxAtStartPar
Alignment timescale (def=None)

\end{fulllineitems}

\index{evolve.py command line option@\spxentry{evolve.py command line option}!\sphinxhyphen{}beammodel@\spxentry{\sphinxhyphen{}beammodel}}\index{\sphinxhyphen{}beammodel@\spxentry{\sphinxhyphen{}beammodel}!evolve.py command line option@\spxentry{evolve.py command line option}}

\begin{fulllineitems}
\phantomsection\label{\detokenize{cmd-line-docs:cmdoption-evolve.py-beammodel}}
\pysigstartsignatures
\pysigline
{\sphinxbfcode{\sphinxupquote{\sphinxhyphen{}beammodel}}\sphinxcode{\sphinxupquote{~\textless{}model\textgreater{}}}}
\pysigstopsignatures
\sphinxAtStartPar
Beaming model to use (def=tm98)

\sphinxAtStartPar
Supported: ‘tm98’, ‘none’, ‘const’, ‘wj08’

\end{fulllineitems}

\index{evolve.py command line option@\spxentry{evolve.py command line option}!\sphinxhyphen{}w@\spxentry{\sphinxhyphen{}w}}\index{\sphinxhyphen{}w@\spxentry{\sphinxhyphen{}w}!evolve.py command line option@\spxentry{evolve.py command line option}}

\begin{fulllineitems}
\phantomsection\label{\detokenize{cmd-line-docs:cmdoption-evolve.py-w}}
\pysigstartsignatures
\pysigline
{\sphinxbfcode{\sphinxupquote{\sphinxhyphen{}w}}\sphinxcode{\sphinxupquote{~\textless{}width\textgreater{}}}}
\pysigstopsignatures
\sphinxAtStartPar
Pulse width to use when generating pulsars, percent (def=5.0)

\end{fulllineitems}

\index{evolve.py command line option@\spxentry{evolve.py command line option}!\sphinxhyphen{}wmod@\spxentry{\sphinxhyphen{}wmod}}\index{\sphinxhyphen{}wmod@\spxentry{\sphinxhyphen{}wmod}!evolve.py command line option@\spxentry{evolve.py command line option}}

\begin{fulllineitems}
\phantomsection\label{\detokenize{cmd-line-docs:cmdoption-evolve.py-wmod}}
\pysigstartsignatures
\pysigline
{\sphinxbfcode{\sphinxupquote{\sphinxhyphen{}wmod}}\sphinxcode{\sphinxupquote{~\textless{}width~model\textgreater{}}}}
\pysigstopsignatures
\sphinxAtStartPar
Pulse width model to use

\sphinxAtStartPar
Supported: None, ‘kj07’

\end{fulllineitems}

\index{evolve.py command line option@\spxentry{evolve.py command line option}!\sphinxhyphen{}sc@\spxentry{\sphinxhyphen{}sc}}\index{\sphinxhyphen{}sc@\spxentry{\sphinxhyphen{}sc}!evolve.py command line option@\spxentry{evolve.py command line option}}

\begin{fulllineitems}
\phantomsection\label{\detokenize{cmd-line-docs:cmdoption-evolve.py-sc}}
\pysigstartsignatures
\pysigline
{\sphinxbfcode{\sphinxupquote{\sphinxhyphen{}sc}}\sphinxcode{\sphinxupquote{~\textless{}scatter~index\textgreater{}}}}
\pysigstopsignatures
\sphinxAtStartPar
Spectral index of scattering law to use (def=\sphinxhyphen{}3.86)

\end{fulllineitems}

\index{evolve.py command line option@\spxentry{evolve.py command line option}!\sphinxhyphen{}eff@\spxentry{\sphinxhyphen{}eff}}\index{\sphinxhyphen{}eff@\spxentry{\sphinxhyphen{}eff}!evolve.py command line option@\spxentry{evolve.py command line option}}

\begin{fulllineitems}
\phantomsection\label{\detokenize{cmd-line-docs:cmdoption-evolve.py-eff}}
\pysigstartsignatures
\pysigline
{\sphinxbfcode{\sphinxupquote{\sphinxhyphen{}eff}}\sphinxcode{\sphinxupquote{~\textless{}efficiency~cutoff\textgreater{}}}}
\pysigstopsignatures
\sphinxAtStartPar
Efficiency cutoff value (def=None)

\end{fulllineitems}

\index{evolve.py command line option@\spxentry{evolve.py command line option}!\sphinxhyphen{}z@\spxentry{\sphinxhyphen{}z}}\index{\sphinxhyphen{}z@\spxentry{\sphinxhyphen{}z}!evolve.py command line option@\spxentry{evolve.py command line option}}

\begin{fulllineitems}
\phantomsection\label{\detokenize{cmd-line-docs:cmdoption-evolve.py-z}}
\pysigstartsignatures
\pysigline
{\sphinxbfcode{\sphinxupquote{\sphinxhyphen{}z}}\sphinxcode{\sphinxupquote{~\textless{}scale~height\textgreater{}}}}
\pysigstopsignatures
\sphinxAtStartPar
Exponential z\sphinxhyphen{}scale height for the population, in kpc (def=0.05)

\end{fulllineitems}

\index{evolve.py command line option@\spxentry{evolve.py command line option}!\sphinxhyphen{}bi@\spxentry{\sphinxhyphen{}bi}}\index{\sphinxhyphen{}bi@\spxentry{\sphinxhyphen{}bi}!evolve.py command line option@\spxentry{evolve.py command line option}}

\begin{fulllineitems}
\phantomsection\label{\detokenize{cmd-line-docs:cmdoption-evolve.py-bi}}
\pysigstartsignatures
\pysigline
{\sphinxbfcode{\sphinxupquote{\sphinxhyphen{}bi}}\sphinxcode{\sphinxupquote{~\textless{}braking~index\textgreater{}}}}
\pysigstopsignatures
\sphinxAtStartPar
Braking index value to use (def=0 = model default)

\end{fulllineitems}

\index{evolve.py command line option@\spxentry{evolve.py command line option}!\sphinxhyphen{}\sphinxhyphen{}nostdout@\spxentry{\sphinxhyphen{}\sphinxhyphen{}nostdout}}\index{\sphinxhyphen{}\sphinxhyphen{}nostdout@\spxentry{\sphinxhyphen{}\sphinxhyphen{}nostdout}!evolve.py command line option@\spxentry{evolve.py command line option}}

\begin{fulllineitems}
\phantomsection\label{\detokenize{cmd-line-docs:cmdoption-evolve.py-nostdout}}
\pysigstartsignatures
\pysigline
{\sphinxbfcode{\sphinxupquote{\sphinxhyphen{}\sphinxhyphen{}nostdout}}\sphinxcode{\sphinxupquote{}}}
\pysigstopsignatures
\sphinxAtStartPar
Turn off writing to stdout.

\end{fulllineitems}

\index{evolve.py command line option@\spxentry{evolve.py command line option}!\sphinxhyphen{}\sphinxhyphen{}nodeathline@\spxentry{\sphinxhyphen{}\sphinxhyphen{}nodeathline}}\index{\sphinxhyphen{}\sphinxhyphen{}nodeathline@\spxentry{\sphinxhyphen{}\sphinxhyphen{}nodeathline}!evolve.py command line option@\spxentry{evolve.py command line option}}

\begin{fulllineitems}
\phantomsection\label{\detokenize{cmd-line-docs:cmdoption-evolve.py-nodeathline}}
\pysigstartsignatures
\pysigline
{\sphinxbfcode{\sphinxupquote{\sphinxhyphen{}\sphinxhyphen{}nodeathline}}\sphinxcode{\sphinxupquote{}}}
\pysigstopsignatures
\sphinxAtStartPar
Turn off the deathline.

\end{fulllineitems}

\index{evolve.py command line option@\spxentry{evolve.py command line option}!\sphinxhyphen{}\sphinxhyphen{}nospiralarms@\spxentry{\sphinxhyphen{}\sphinxhyphen{}nospiralarms}}\index{\sphinxhyphen{}\sphinxhyphen{}nospiralarms@\spxentry{\sphinxhyphen{}\sphinxhyphen{}nospiralarms}!evolve.py command line option@\spxentry{evolve.py command line option}}

\begin{fulllineitems}
\phantomsection\label{\detokenize{cmd-line-docs:cmdoption-evolve.py-nospiralarms}}
\pysigstartsignatures
\pysigline
{\sphinxbfcode{\sphinxupquote{\sphinxhyphen{}\sphinxhyphen{}nospiralarms}}\sphinxcode{\sphinxupquote{}}}
\pysigstopsignatures
\sphinxAtStartPar
Turn off spiral arms galactic distribution.

\end{fulllineitems}

\index{evolve.py command line option@\spxentry{evolve.py command line option}!\sphinxhyphen{}\sphinxhyphen{}keepdead@\spxentry{\sphinxhyphen{}\sphinxhyphen{}keepdead}}\index{\sphinxhyphen{}\sphinxhyphen{}keepdead@\spxentry{\sphinxhyphen{}\sphinxhyphen{}keepdead}!evolve.py command line option@\spxentry{evolve.py command line option}}

\begin{fulllineitems}
\phantomsection\label{\detokenize{cmd-line-docs:cmdoption-evolve.py-keepdead}}
\pysigstartsignatures
\pysigline
{\sphinxbfcode{\sphinxupquote{\sphinxhyphen{}\sphinxhyphen{}keepdead}}\sphinxcode{\sphinxupquote{}}}
\pysigstopsignatures
\sphinxAtStartPar
Keep dead pulsars in the population model.

\end{fulllineitems}



\section{view.py}
\label{\detokenize{cmd-line-docs:view-py}}\label{\detokenize{cmd-line-docs:view-doc}}\label{\detokenize{cmd-line-docs:view-docs}}\index{view.py command line option@\spxentry{view.py command line option}!\sphinxhyphen{}f@\spxentry{\sphinxhyphen{}f}}\index{\sphinxhyphen{}f@\spxentry{\sphinxhyphen{}f}!view.py command line option@\spxentry{view.py command line option}}

\begin{fulllineitems}
\phantomsection\label{\detokenize{cmd-line-docs:cmdoption-view.py-f}}
\pysigstartsignatures
\pysigline
{\sphinxbfcode{\sphinxupquote{\sphinxhyphen{}f}}\sphinxcode{\sphinxupquote{~\textless{}input~model\textgreater{}}}}
\pysigstopsignatures
\sphinxAtStartPar
Select the population model to view (default=populate.model)

\end{fulllineitems}

\index{view.py command line option@\spxentry{view.py command line option}!\sphinxhyphen{}p@\spxentry{\sphinxhyphen{}p}}\index{\sphinxhyphen{}p@\spxentry{\sphinxhyphen{}p}!view.py command line option@\spxentry{view.py command line option}}

\begin{fulllineitems}
\phantomsection\label{\detokenize{cmd-line-docs:cmdoption-view.py-p}}
\pysigstartsignatures
\pysigline
{\sphinxbfcode{\sphinxupquote{\sphinxhyphen{}p}}\sphinxcode{\sphinxupquote{~\textless{}param~name\textgreater{}}}}
\pysigstopsignatures
\sphinxAtStartPar
Select the population parameter to plot

\sphinxAtStartPar
Supported: ‘period’, ‘dm’, ‘gl’, ‘gb’, ‘lum’, ‘alpha’, ‘r0’, ‘rho’, ‘width’, ‘spindex’, ‘scindex’, ‘dist’

\end{fulllineitems}

\index{view.py command line option@\spxentry{view.py command line option}!\sphinxhyphen{}\sphinxhyphen{}logx@\spxentry{\sphinxhyphen{}\sphinxhyphen{}logx}}\index{\sphinxhyphen{}\sphinxhyphen{}logx@\spxentry{\sphinxhyphen{}\sphinxhyphen{}logx}!view.py command line option@\spxentry{view.py command line option}}

\begin{fulllineitems}
\phantomsection\label{\detokenize{cmd-line-docs:cmdoption-view.py-logx}}
\pysigstartsignatures
\pysigline
{\sphinxbfcode{\sphinxupquote{\sphinxhyphen{}\sphinxhyphen{}logx}}\sphinxcode{\sphinxupquote{}}}
\pysigstopsignatures
\sphinxAtStartPar
Plot log10 of the values

\end{fulllineitems}



\section{visualize.py}
\label{\detokenize{cmd-line-docs:visualize-py}}\label{\detokenize{cmd-line-docs:visualize-docs}}\index{visualize.py command line option@\spxentry{visualize.py command line option}!\sphinxhyphen{}f@\spxentry{\sphinxhyphen{}f}}\index{\sphinxhyphen{}f@\spxentry{\sphinxhyphen{}f}!visualize.py command line option@\spxentry{visualize.py command line option}}

\begin{fulllineitems}
\phantomsection\label{\detokenize{cmd-line-docs:cmdoption-visualize.py-f}}
\pysigstartsignatures
\pysigline
{\sphinxbfcode{\sphinxupquote{\sphinxhyphen{}f}}\sphinxcode{\sphinxupquote{~\textless{}model\textgreater{}}}}
\pysigstopsignatures
\sphinxAtStartPar
Select a population model to plot (default = populate.model)

\end{fulllineitems}

\index{visualize.py command line option@\spxentry{visualize.py command line option}!\sphinxhyphen{}frac@\spxentry{\sphinxhyphen{}frac}}\index{\sphinxhyphen{}frac@\spxentry{\sphinxhyphen{}frac}!visualize.py command line option@\spxentry{visualize.py command line option}}

\begin{fulllineitems}
\phantomsection\label{\detokenize{cmd-line-docs:cmdoption-visualize.py-frac}}
\pysigstartsignatures
\pysigline
{\sphinxbfcode{\sphinxupquote{\sphinxhyphen{}frac}}\sphinxcode{\sphinxupquote{~\textless{}F\textgreater{}}}}
\pysigstopsignatures
\sphinxAtStartPar
Plot a fraction \textless{}F\textgreater{} of the population for speed gains

\end{fulllineitems}


\sphinxstepscope


\chapter{Tutorial \sphinxhyphen{} Basic Usage}
\label{\detokenize{tutorial-basic:tutorial-basic-usage}}\label{\detokenize{tutorial-basic:tutorial-basic}}\label{\detokenize{tutorial-basic::doc}}
\sphinxAtStartPar
This page will outline a very simple pipeline for basic population
simulations with PsrPopPy.

\sphinxAtStartPar
The following describe the simple command line workflow with the psrpoppysuper scripts.


\section{Generate Population Model}
\label{\detokenize{tutorial-basic:generate-population-model}}\label{\detokenize{tutorial-basic:generate-population}}
\sphinxAtStartPar
Population models are generated using the \sphinxcode{\sphinxupquote{populate}} script (installed
automatically with pip). A common use would be to generate a population of
normal pulsars using the Parkes Multibeam Pulsar Survey as a basis. This
survey detected 1038 pulsars (at last count). Using default radial distribution,
period and luminosity models, we can generate such a population using the command:

\begin{sphinxVerbatim}[commandchars=\\\{\}]
\PYG{n}{populate} \PYG{o}{\PYGZhy{}}\PYG{n}{n} \PYG{l+m+mi}{1038} \PYG{o}{\PYGZhy{}}\PYG{n}{surveys} \PYG{n}{PMSURV}
\end{sphinxVerbatim}

\sphinxAtStartPar
This will then run for a few minutes, until the model PMSURV survey
has detected 1038 pulsars. The file \sphinxcode{\sphinxupquote{populate.model}} will be produced
by default, which is a \sphinxhref{http://docs.python.org/library/pickle.html}{pickled}
population object.


\section{Evolve a Population Model}
\label{\detokenize{tutorial-basic:evolve-a-population-model}}\label{\detokenize{tutorial-basic:evolve-population}}
\sphinxAtStartPar
The \sphinxcode{\sphinxupquote{evolve}} script takes a snapshot\sphinxhyphen{}style population model and evolves
it according to the Ridley \& Lorimer pulsar evolution method. This is
useful when you want to generate a more realistic, age\sphinxhyphen{}dependent
population before passing it on to a survey simulation.

\sphinxAtStartPar
A basic usage example is:

\begin{sphinxVerbatim}[commandchars=\\\{\}]
\PYG{n}{evolve} \PYG{o}{\PYGZhy{}}\PYG{n}{n} \PYG{l+m+mi}{1000} \PYG{o}{\PYGZhy{}}\PYG{n}{surveys} \PYG{n}{PMSURV} \PYG{o}{\PYGZhy{}}\PYG{n}{dm} \PYG{n}{ne2025} \PYG{o}{\PYGZhy{}}\PYG{n}{o} \PYG{n}{evolve}\PYG{o}{.}\PYG{n}{model}
\end{sphinxVerbatim}

\sphinxAtStartPar
This command will generate an evolved population of pulsars
using the NE2025 electron density model. The output file
\sphinxcode{\sphinxupquote{evolve.model}} can then be used by \sphinxcode{\sphinxupquote{dosurvey}} or visualised with
\sphinxcode{\sphinxupquote{view.py}}.

\sphinxAtStartPar
The evolve script also supports the same survey and population controls as
\sphinxcode{\sphinxupquote{populate}}, including pulse width, beaming, spin\sphinxhyphen{}down, and electron
model options. For example, to use a fixed 5\% pulse width and the WJ08
alignment model, run:

\begin{sphinxVerbatim}[commandchars=\\\{\}]
\PYG{n}{evolve} \PYG{o}{\PYGZhy{}}\PYG{n}{n} \PYG{l+m+mi}{1000} \PYG{o}{\PYGZhy{}}\PYG{n}{surveys} \PYG{n}{PMSURV} \PYG{o}{\PYGZhy{}}\PYG{n}{w} \PYG{l+m+mf}{5.0} \PYG{o}{\PYGZhy{}}\PYG{n}{alignmodel} \PYG{n}{wj08} \PYG{o}{\PYGZhy{}}\PYG{n}{o} \PYG{n}{evolve}\PYG{o}{.}\PYG{n}{model}
\end{sphinxVerbatim}

\sphinxAtStartPar
If you only want to generate the evolved population without checking a
survey, omit the \sphinxcode{\sphinxupquote{\sphinxhyphen{}surveys}} option and use the output file directly:
\begin{quote}

\sphinxAtStartPar
evolve \sphinxhyphen{}n 1000 \sphinxhyphen{}dm ne2025 \sphinxhyphen{}o evolve.model
\end{quote}


\section{Simulate a Pulsar Survey}
\label{\detokenize{tutorial-basic:simulate-a-pulsar-survey}}\label{\detokenize{tutorial-basic:simulate-survey}}
\sphinxAtStartPar
Once you’ve generated a pulsar population model, the \sphinxcode{\sphinxupquote{dosurvey}} script
can be used to run the model through a past, present or future, pulsar
survey (as specified in files in the \sphinxcode{\sphinxupquote{survey}} directory shipped with the code
.. \DUrole{xref}{\DUrole{std}{\DUrole{std-ref}{\_model\sphinxhyphen{}survey\sphinxhyphen{}files}}}).

\sphinxAtStartPar
For example, say we want to take the population model we just created,
\sphinxcode{\sphinxupquote{pop\_ne2025.model}}, and estimate from this how many pulsars would be detected
in a putative LOFAR pulsar survey. This can be simply done using:

\begin{sphinxVerbatim}[commandchars=\\\{\}]
\PYG{n}{dosurvey} \PYG{o}{\PYGZhy{}}\PYG{n}{f} \PYG{n}{pop\PYGZus{}ne2025}\PYG{o}{.}\PYG{n}{model} \PYG{o}{\PYGZhy{}}\PYG{n}{surveys} \PYG{n}{LOFAR}
\end{sphinxVerbatim}


\section{Using PsrPopPySuper as a Python Module}
\label{\detokenize{tutorial-basic:using-psrpoppysuper-as-a-python-module}}
\sphinxAtStartPar
The PsrPopPySuper package can also be used as a Python module. The
following example shows how to generate a population model, evolve it, and
then run it through a survey, all in one script:

\begin{sphinxVerbatim}[commandchars=\\\{\}]
\PYG{k+kn}{from}\PYG{+w}{ }\PYG{n+nn}{psrpoppysuper}\PYG{+w}{ }\PYG{k+kn}{import} \PYG{n}{evolve}\PYG{p}{,}\PYG{n}{dosurvey}

\PYG{n}{evolution} \PYG{o}{=} \PYG{n}{evolve}\PYG{o}{.}\PYG{n}{generate}\PYG{p}{(}\PYG{n}{ngen}\PYG{o}{=}\PYG{l+m+mi}{1269}\PYG{p}{,}
          \PYG{c+c1}{\PYGZsh{} surveyList=[\PYGZsq{}PMPS\PYGZus{}NEW\PYGZsq{}],}
          \PYG{n}{age\PYGZus{}max}\PYG{o}{=}\PYG{l+m+mf}{1.0E+9}\PYG{p}{,}
          \PYG{n}{birthVModel}\PYG{o}{=}\PYG{l+s+s1}{\PYGZsq{}}\PYG{l+s+s1}{exp}\PYG{l+s+s1}{\PYGZsq{}}\PYG{p}{,}
          \PYG{n}{birthVPars}\PYG{o}{=}\PYG{p}{[}\PYG{l+m+mi}{0}\PYG{p}{,}\PYG{l+m+mi}{380}\PYG{p}{]}\PYG{p}{,}
          \PYG{n}{electronModel}\PYG{o}{=}\PYG{l+s+s1}{\PYGZsq{}}\PYG{l+s+s1}{ne2001}\PYG{l+s+s1}{\PYGZsq{}}\PYG{p}{,}
          \PYG{n}{bFieldPars}\PYG{o}{=}\PYG{p}{[}\PYG{l+m+mf}{12.35}\PYG{p}{,}\PYG{l+m+mf}{0.55}\PYG{p}{]}\PYG{p}{,}
          \PYG{n}{lumDistType}\PYG{o}{=}\PYG{l+s+s1}{\PYGZsq{}}\PYG{l+s+s1}{fk06}\PYG{l+s+s1}{\PYGZsq{}}\PYG{p}{,}
          \PYG{n}{lumDistPars}\PYG{o}{=}\PYG{p}{[}\PYG{o}{\PYGZhy{}}\PYG{l+m+mf}{1.5}\PYG{p}{,}\PYG{l+m+mf}{0.5}\PYG{p}{]}\PYG{p}{,}
          \PYG{n}{pDistPars}\PYG{o}{=}\PYG{p}{[}\PYG{l+m+mf}{0.3}\PYG{p}{,}\PYG{l+m+mf}{0.15}\PYG{p}{]}\PYG{p}{,}
          \PYG{n}{braking\PYGZus{}index}\PYG{o}{=}\PYG{l+m+mi}{3}
\PYG{p}{)}

\PYG{n}{survey} \PYG{o}{=} \PYG{n}{dosurvey}\PYG{o}{.}\PYG{n}{run}\PYG{p}{(}\PYG{n}{evolution}\PYG{p}{,}\PYG{n}{surveyList}\PYG{o}{=}\PYG{p}{[}\PYG{l+s+s1}{\PYGZsq{}}\PYG{l+s+s1}{PMPS\PYGZus{}NEW}\PYG{l+s+s1}{\PYGZsq{}}\PYG{p}{]}\PYG{p}{)}

\PYG{n}{dosurvey}\PYG{o}{.}\PYG{n}{write}\PYG{p}{(}\PYG{n}{survey}\PYG{p}{,}\PYG{n}{asc}\PYG{o}{=}\PYG{k+kc}{True}\PYG{p}{,}\PYG{n}{extension}\PYG{o}{=}\PYG{l+s+sa}{f}\PYG{l+s+s1}{\PYGZsq{}}\PYG{l+s+si}{\PYGZob{}}\PYG{n}{i}\PYG{l+s+si}{\PYGZcb{}}\PYG{l+s+s1}{.results}\PYG{l+s+s1}{\PYGZsq{}}\PYG{p}{)}
\end{sphinxVerbatim}

\sphinxAtStartPar
These runs the PsrpopPySuper functions directly, and the output is a population object which can be
used for further analysis or visualisation. The \sphinxcode{\sphinxupquote{write}} function can be used to save the results to a file,
with the option to save as an ascii file for easy inspection.

\sphinxstepscope


\chapter{Module\sphinxhyphen{}level Documentation}
\label{\detokenize{module-level-docs:module-level-documentation}}\label{\detokenize{module-level-docs:module-level-docs}}\label{\detokenize{module-level-docs::doc}}

\section{{\hyperref[\detokenize{module-level-docs:module-pulsar}]{\sphinxcrossref{\sphinxstyleliteralintitle{\sphinxupquote{pulsar}}}}} \textendash{} Creates/stores a pulsar object}
\label{\detokenize{module-level-docs:module-pulsar}}\label{\detokenize{module-level-docs:pulsar-creates-stores-a-pulsar-object}}\label{\detokenize{module-level-docs:pulsar-module-docs}}\index{module@\spxentry{module}!pulsar@\spxentry{pulsar}}\index{pulsar@\spxentry{pulsar}!module@\spxentry{module}}\index{Pulsar (class in pulsar)@\spxentry{Pulsar}\spxextra{class in pulsar}}

\begin{fulllineitems}
\phantomsection\label{\detokenize{module-level-docs:pulsar.Pulsar}}
\pysigstartsignatures
\pysigline
{\sphinxbfcode{\sphinxupquote{\DUrole{k}{class}\DUrole{w}{ }}}\sphinxcode{\sphinxupquote{pulsar.}}\sphinxbfcode{\sphinxupquote{Pulsar}}}
\pysigstopsignatures
\end{fulllineitems}

\index{\_\_init\_\_() (pulsar.Pulsar method)@\spxentry{\_\_init\_\_()}\spxextra{pulsar.Pulsar method}}

\begin{fulllineitems}
\phantomsection\label{\detokenize{module-level-docs:pulsar.Pulsar.__init__}}
\pysigstartsignatures
\pysiglinewithargsret
{\sphinxcode{\sphinxupquote{Pulsar.}}\sphinxbfcode{\sphinxupquote{\_\_init\_\_}}}
{\sphinxoptional{\sphinxparam{\DUrole{n}{period}}\sphinxparamcomma \sphinxparam{\DUrole{n}{dm}}\sphinxparamcomma \sphinxparam{\DUrole{n}{gl}}\sphinxparamcomma \sphinxparam{\DUrole{n}{gb}}\sphinxparamcomma \sphinxparam{\DUrole{n}{galCoords}}\sphinxparamcomma \sphinxparam{\DUrole{n}{r0}}\sphinxparamcomma \sphinxparam{\DUrole{n}{dtrue}}\sphinxparamcomma \sphinxparam{\DUrole{n}{lum\_1400}}\sphinxparamcomma \sphinxparam{\DUrole{n}{spindex}}\sphinxparamcomma \sphinxparam{\DUrole{n}{alpha}}\sphinxparamcomma \sphinxparam{\DUrole{n}{rho}}\sphinxparamcomma \sphinxparam{\DUrole{n}{width\_degree}}\sphinxparamcomma \sphinxparam{\DUrole{n}{snr}}\sphinxparamcomma \sphinxparam{\DUrole{n}{beaming}}\sphinxparamcomma \sphinxparam{\DUrole{n}{scindex}}\sphinxparamcomma \sphinxparam{\DUrole{n}{gpsFlag}}\sphinxparamcomma \sphinxparam{\DUrole{n}{gpsA}}\sphinxparamcomma \sphinxparam{\DUrole{n}{brokenFlag}}\sphinxparamcomma \sphinxparam{\DUrole{n}{brokenSI}}}}
{}
\pysigstopsignatures
\sphinxAtStartPar
Initialise the pulsar object

\end{fulllineitems}

\index{s\_1400() (pulsar.Pulsar method)@\spxentry{s\_1400()}\spxextra{pulsar.Pulsar method}}

\begin{fulllineitems}
\phantomsection\label{\detokenize{module-level-docs:pulsar.Pulsar.s_1400}}
\pysigstartsignatures
\pysiglinewithargsret
{\sphinxcode{\sphinxupquote{Pulsar.}}\sphinxbfcode{\sphinxupquote{s\_1400}}}
{}
{}
\pysigstopsignatures
\sphinxAtStartPar
Returns the flux at 1400 MHz, calculated as

\sphinxAtStartPar
\(S_{\rm{1400}} = \frac{L_{\rm{1400}}}{D_{\rm{true}}^2}\)

\end{fulllineitems}

\index{width\_ms() (pulsar.Pulsar method)@\spxentry{width\_ms()}\spxextra{pulsar.Pulsar method}}

\begin{fulllineitems}
\phantomsection\label{\detokenize{module-level-docs:pulsar.Pulsar.width_ms}}
\pysigstartsignatures
\pysiglinewithargsret
{\sphinxcode{\sphinxupquote{Pulsar.}}\sphinxbfcode{\sphinxupquote{width\_ms}}}
{}
{}
\pysigstopsignatures
\sphinxAtStartPar
Returns the pulse width in milliseconds, calculated as

\sphinxAtStartPar
\(W_{\rm{ms}} = P_{\rm{ms}} \times W_{\rm{degree}} / 360\)

\end{fulllineitems}



\section{{\hyperref[\detokenize{module-level-docs:module-population}]{\sphinxcrossref{\sphinxstyleliteralintitle{\sphinxupquote{population}}}}} \textendash{} Creates/stores a population}
\label{\detokenize{module-level-docs:module-population}}\label{\detokenize{module-level-docs:population-creates-stores-a-population}}\label{\detokenize{module-level-docs:population-module-docs}}\index{module@\spxentry{module}!population@\spxentry{population}}\index{population@\spxentry{population}!module@\spxentry{module}}\index{Population (class in population)@\spxentry{Population}\spxextra{class in population}}

\begin{fulllineitems}
\phantomsection\label{\detokenize{module-level-docs:population.Population}}
\pysigstartsignatures
\pysigline
{\sphinxbfcode{\sphinxupquote{\DUrole{k}{class}\DUrole{w}{ }}}\sphinxcode{\sphinxupquote{population.}}\sphinxbfcode{\sphinxupquote{Population}}}
\pysigstopsignatures
\end{fulllineitems}

\index{\_\_init\_\_() (population.Population method)@\spxentry{\_\_init\_\_()}\spxextra{population.Population method}}

\begin{fulllineitems}
\phantomsection\label{\detokenize{module-level-docs:population.Population.__init__}}
\pysigstartsignatures
\pysiglinewithargsret
{\sphinxcode{\sphinxupquote{Population.}}\sphinxbfcode{\sphinxupquote{\_\_init\_\_}}}
{\sphinxoptional{\sphinxparam{\DUrole{n}{pDistType}}\sphinxparamcomma \sphinxparam{\DUrole{n}{radialDistType}}\sphinxparamcomma \sphinxparam{\DUrole{n}{lumDistType}}\sphinxparamcomma \sphinxparam{\DUrole{n}{pmean}}\sphinxparamcomma \sphinxparam{\DUrole{n}{psigma}}\sphinxparamcomma \sphinxparam{\DUrole{n}{simean}}\sphinxparamcomma \sphinxparam{\DUrole{n}{sisigma}}\sphinxparamcomma \sphinxparam{\DUrole{n}{lummean}}\sphinxparamcomma \sphinxparam{\DUrole{n}{lumsigma}}\sphinxparamcomma \sphinxparam{\DUrole{n}{zscale}}\sphinxparamcomma \sphinxparam{\DUrole{n}{electronModel}}\sphinxparamcomma \sphinxparam{\DUrole{n}{gpsFrac}}\sphinxparamcomma \sphinxparam{\DUrole{n}{gpsA}}\sphinxparamcomma \sphinxparam{\DUrole{n}{brokenFrac}}\sphinxparamcomma \sphinxparam{\DUrole{n}{brokenSI}}\sphinxparamcomma \sphinxparam{\DUrole{n}{ref\_freq}}}}
{}
\pysigstopsignatures
\sphinxAtStartPar
Initialise the population object

\end{fulllineitems}

\index{\_\_str\_\_() (population.Population method)@\spxentry{\_\_str\_\_()}\spxextra{population.Population method}}

\begin{fulllineitems}
\phantomsection\label{\detokenize{module-level-docs:population.Population.__str__}}
\pysigstartsignatures
\pysiglinewithargsret
{\sphinxcode{\sphinxupquote{Population.}}\sphinxbfcode{\sphinxupquote{\_\_str\_\_}}}
{}
{}
\pysigstopsignatures
\sphinxAtStartPar
Defines how the operation \sphinxcode{\sphinxupquote{print Population}} is performed

\end{fulllineitems}

\index{size() (population.Population method)@\spxentry{size()}\spxextra{population.Population method}}

\begin{fulllineitems}
\phantomsection\label{\detokenize{module-level-docs:population.Population.size}}
\pysigstartsignatures
\pysiglinewithargsret
{\sphinxcode{\sphinxupquote{Population.}}\sphinxbfcode{\sphinxupquote{size}}}
{}
{}
\pysigstopsignatures
\sphinxAtStartPar
Returns the number of pulsars in the population object

\end{fulllineitems}

\index{join() (population.Population method)@\spxentry{join()}\spxextra{population.Population method}}

\begin{fulllineitems}
\phantomsection\label{\detokenize{module-level-docs:population.Population.join}}
\pysigstartsignatures
\pysiglinewithargsret
{\sphinxcode{\sphinxupquote{Population.}}\sphinxbfcode{\sphinxupquote{join}}}
{\sphinxparam{\DUrole{n}{poplist}}}
{}
\pysigstopsignatures
\sphinxAtStartPar
Joins each of the populations in list \sphinxcode{\sphinxupquote{poplist}} to the current population

\end{fulllineitems}

\index{write() (population.Population method)@\spxentry{write()}\spxextra{population.Population method}}

\begin{fulllineitems}
\phantomsection\label{\detokenize{module-level-docs:population.Population.write}}
\pysigstartsignatures
\pysiglinewithargsret
{\sphinxcode{\sphinxupquote{Population.}}\sphinxbfcode{\sphinxupquote{write}}}
{\sphinxparam{\DUrole{n}{outf}}}
{}
\pysigstopsignatures
\sphinxAtStartPar
Uses cPickle to dump the population to file \sphinxcode{\sphinxupquote{outf}}

\end{fulllineitems}

\index{write\_asc() (population.Population method)@\spxentry{write\_asc()}\spxextra{population.Population method}}

\begin{fulllineitems}
\phantomsection\label{\detokenize{module-level-docs:population.Population.write_asc}}
\pysigstartsignatures
\pysiglinewithargsret
{\sphinxcode{\sphinxupquote{Population.}}\sphinxbfcode{\sphinxupquote{write\_asc}}}
{\sphinxparam{\DUrole{n}{outf}}}
{}
\pysigstopsignatures
\sphinxAtStartPar
Writes the population to an ascii file in the old psrpop way

\end{fulllineitems}



\section{{\hyperref[\detokenize{module-level-docs:module-survey}]{\sphinxcrossref{\sphinxstyleliteralintitle{\sphinxupquote{survey}}}}} \textendash{} Read a survey file into a survey object}
\label{\detokenize{module-level-docs:module-survey}}\label{\detokenize{module-level-docs:survey-read-a-survey-file-into-a-survey-object}}\label{\detokenize{module-level-docs:survey-module-docs}}\index{module@\spxentry{module}!survey@\spxentry{survey}}\index{survey@\spxentry{survey}!module@\spxentry{module}}\index{Survey (class in survey)@\spxentry{Survey}\spxextra{class in survey}}

\begin{fulllineitems}
\phantomsection\label{\detokenize{module-level-docs:survey.Survey}}
\pysigstartsignatures
\pysigline
{\sphinxbfcode{\sphinxupquote{\DUrole{k}{class}\DUrole{w}{ }}}\sphinxcode{\sphinxupquote{survey.}}\sphinxbfcode{\sphinxupquote{Survey}}}
\pysigstopsignatures
\end{fulllineitems}

\index{\_\_init\_\_() (survey.Survey method)@\spxentry{\_\_init\_\_()}\spxextra{survey.Survey method}}

\begin{fulllineitems}
\phantomsection\label{\detokenize{module-level-docs:survey.Survey.__init__}}
\pysigstartsignatures
\pysiglinewithargsret
{\sphinxcode{\sphinxupquote{Survey.}}\sphinxbfcode{\sphinxupquote{\_\_init\_\_}}}
{\sphinxparam{\DUrole{n}{surveyName}}}
{}
\pysigstopsignatures
\sphinxAtStartPar
Read in a (correctly formatted!) survey file

\end{fulllineitems}

\index{\_\_str\_\_() (survey.Survey method)@\spxentry{\_\_str\_\_()}\spxextra{survey.Survey method}}

\begin{fulllineitems}
\phantomsection\label{\detokenize{module-level-docs:survey.Survey.__str__}}
\pysigstartsignatures
\pysiglinewithargsret
{\sphinxcode{\sphinxupquote{Survey.}}\sphinxbfcode{\sphinxupquote{\_\_str\_\_}}}
{}
{}
\pysigstopsignatures
\sphinxAtStartPar
Define how to perform \sphinxcode{\sphinxupquote{print Survey}}

\end{fulllineitems}

\index{nchans() (survey.Survey method)@\spxentry{nchans()}\spxextra{survey.Survey method}}

\begin{fulllineitems}
\phantomsection\label{\detokenize{module-level-docs:survey.Survey.nchans}}
\pysigstartsignatures
\pysiglinewithargsret
{\sphinxcode{\sphinxupquote{Survey.}}\sphinxbfcode{\sphinxupquote{nchans}}}
{}
{}
\pysigstopsignatures
\sphinxAtStartPar
Returns the number of channels, calculated as

\sphinxAtStartPar
\(n_{\rm{chans}} = \frac{ {\rm{BW}}_{\rm{total}} }{ {\rm{BM}}_{\rm{chan}} }\)

\end{fulllineitems}

\index{inRegion() (survey.Survey method)@\spxentry{inRegion()}\spxextra{survey.Survey method}}

\begin{fulllineitems}
\phantomsection\label{\detokenize{module-level-docs:survey.Survey.inRegion}}
\pysigstartsignatures
\pysiglinewithargsret
{\sphinxcode{\sphinxupquote{Survey.}}\sphinxbfcode{\sphinxupquote{inRegion}}}
{\sphinxparam{\DUrole{n}{pulsar}}}
{}
\pysigstopsignatures
\sphinxAtStartPar
Determines if {\hyperref[\detokenize{module-level-docs:pulsar.Pulsar}]{\sphinxcrossref{\sphinxcode{\sphinxupquote{Pulsar}}}}} is inside survey region. Returns True or False accordingly

\end{fulllineitems}

\index{inPointing() (survey.Survey method)@\spxentry{inPointing()}\spxextra{survey.Survey method}}

\begin{fulllineitems}
\phantomsection\label{\detokenize{module-level-docs:survey.Survey.inPointing}}
\pysigstartsignatures
\pysiglinewithargsret
{\sphinxcode{\sphinxupquote{Survey.}}\sphinxbfcode{\sphinxupquote{inPointing}}}
{\sphinxparam{\DUrole{n}{pulsar}}}
{}
\pysigstopsignatures
\sphinxAtStartPar
Determines if {\hyperref[\detokenize{module-level-docs:pulsar.Pulsar}]{\sphinxcrossref{\sphinxcode{\sphinxupquote{Pulsar}}}}} is inside one of the survey’s {\hyperref[\detokenize{module-level-docs:survey.Pointing}]{\sphinxcrossref{\sphinxcode{\sphinxupquote{pointings}}}}}. Returns the offset from beam centre to the pulsar.

\end{fulllineitems}

\index{SNRcalc() (survey.Survey method)@\spxentry{SNRcalc()}\spxextra{survey.Survey method}}

\begin{fulllineitems}
\phantomsection\label{\detokenize{module-level-docs:survey.Survey.SNRcalc}}
\pysigstartsignatures
\pysiglinewithargsret
{\sphinxcode{\sphinxupquote{Survey.}}\sphinxbfcode{\sphinxupquote{SNRcalc}}}
{\sphinxparam{\DUrole{n}{pulsar}}\sphinxparamcomma \sphinxparam{\DUrole{n}{pop}}}
{}
\pysigstopsignatures
\sphinxAtStartPar
Calculates the SNR of a {\hyperref[\detokenize{module-level-docs:pulsar.Pulsar}]{\sphinxcrossref{\sphinxcode{\sphinxupquote{Pulsar}}}}} from {\hyperref[\detokenize{module-level-docs:population.Population}]{\sphinxcrossref{\sphinxcode{\sphinxupquote{Population}}}}} \sphinxcode{\sphinxupquote{pop}} in the survey. Returns \sphinxhyphen{}1 if pulse is smeared, and \sphinxhyphen{}2 if pulsar is outside survey region. SNR is calculated (with familiar terms) as

\sphinxAtStartPar
\({\rm{SNR}} = \frac{S_{1400} G \sqrt{n_{\rm{pol}} BW \tau}}{\beta T_{\rm{tot}}} \sqrt{\frac{1-\delta}{\delta}} \eta\)

\sphinxAtStartPar
where

\sphinxAtStartPar
\(\eta = \exp(-2.7727 \times {\rm{offset}}^2 / {\rm{fwhm}}^2)\)

\end{fulllineitems}

\index{Pointing (class in survey)@\spxentry{Pointing}\spxextra{class in survey}}

\begin{fulllineitems}
\phantomsection\label{\detokenize{module-level-docs:survey.Pointing}}
\pysigstartsignatures
\pysigline
{\sphinxbfcode{\sphinxupquote{\DUrole{k}{class}\DUrole{w}{ }}}\sphinxcode{\sphinxupquote{survey.}}\sphinxbfcode{\sphinxupquote{Pointing}}}
\pysigstopsignatures
\end{fulllineitems}

\index{\_\_init\_\_() (survey.Pointing method)@\spxentry{\_\_init\_\_()}\spxextra{survey.Pointing method}}

\begin{fulllineitems}
\phantomsection\label{\detokenize{module-level-docs:survey.Pointing.__init__}}
\pysigstartsignatures
\pysiglinewithargsret
{\sphinxcode{\sphinxupquote{Pointing.}}\sphinxbfcode{\sphinxupquote{\_\_init\_\_}}}
{\sphinxparam{\DUrole{n}{coord1}}\sphinxparamcomma \sphinxparam{\DUrole{n}{coord2}}\sphinxparamcomma \sphinxparam{\DUrole{n}{coordtype}}}
{}
\pysigstopsignatures
\sphinxAtStartPar
Converts (coord1, coord2) into correctly formatted (l, b). Coordtype must be
either \sphinxcode{\sphinxupquote{eq}} or \sphinxcode{\sphinxupquote{gal}}. If \sphinxcode{\sphinxupquote{eq}}, the RA and Dec are converted internally

\end{fulllineitems}



\section{{\hyperref[\detokenize{module-level-docs:module-populate}]{\sphinxcrossref{\sphinxstyleliteralintitle{\sphinxupquote{populate}}}}} \textendash{} Create a population object}
\label{\detokenize{module-level-docs:module-populate}}\label{\detokenize{module-level-docs:populate-create-a-population-object}}\label{\detokenize{module-level-docs:populate-module-docs}}\index{module@\spxentry{module}!populate@\spxentry{populate}}\index{populate@\spxentry{populate}!module@\spxentry{module}}\index{Populate (class in populate)@\spxentry{Populate}\spxextra{class in populate}}

\begin{fulllineitems}
\phantomsection\label{\detokenize{module-level-docs:populate.Populate}}
\pysigstartsignatures
\pysigline
{\sphinxbfcode{\sphinxupquote{\DUrole{k}{class}\DUrole{w}{ }}}\sphinxcode{\sphinxupquote{populate.}}\sphinxbfcode{\sphinxupquote{Populate}}}
\pysigstopsignatures
\end{fulllineitems}

\index{generate() (populate.Populate method)@\spxentry{generate()}\spxextra{populate.Populate method}}

\begin{fulllineitems}
\phantomsection\label{\detokenize{module-level-docs:populate.Populate.generate}}
\pysigstartsignatures
\pysiglinewithargsret
{\sphinxcode{\sphinxupquote{Populate.}}\sphinxbfcode{\sphinxupquote{generate}}}
{\sphinxparam{\DUrole{n}{ngen}}\sphinxoptional{\sphinxparamcomma \sphinxparam{\DUrole{n}{surveyList}}\sphinxparamcomma \sphinxparam{\DUrole{n}{pDistType}}\sphinxparamcomma \sphinxparam{\DUrole{n}{radialDistType}}\sphinxparamcomma \sphinxparam{\DUrole{n}{electronModel}}\sphinxparamcomma \sphinxparam{\DUrole{n}{pDistPars}}\sphinxparamcomma \sphinxparam{\DUrole{n}{siDistPars}}\sphinxparamcomma \sphinxparam{\DUrole{n}{lumDistType}}\sphinxparamcomma \sphinxparam{\DUrole{n}{lumDistPars}}\sphinxparamcomma \sphinxparam{\DUrole{n}{zscale}}\sphinxparamcomma \sphinxparam{\DUrole{n}{duty}}\sphinxparamcomma \sphinxparam{\DUrole{n}{scindex}}\sphinxparamcomma \sphinxparam{\DUrole{n}{gpsArgs}}\sphinxparamcomma \sphinxparam{\DUrole{n}{doubleSpec}}\sphinxparamcomma \sphinxparam{\DUrole{n}{nostdout}}}}
{}
\pysigstopsignatures
\sphinxAtStartPar
The method called by the \sphinxcode{\sphinxupquote{populate.py}} command\sphinxhyphen{}line\sphinxhyphen{}script

\end{fulllineitems}

\index{write() (populate.Populate method)@\spxentry{write()}\spxextra{populate.Populate method}}

\begin{fulllineitems}
\phantomsection\label{\detokenize{module-level-docs:populate.Populate.write}}
\pysigstartsignatures
\pysiglinewithargsret
{\sphinxcode{\sphinxupquote{Populate.}}\sphinxbfcode{\sphinxupquote{write}}}
{\sphinxparam{\DUrole{n}{outf}\DUrole{o}{=}\DUrole{default_value}{populate.model}}}
{}
\pysigstopsignatures
\sphinxAtStartPar
Writes the {\hyperref[\detokenize{module-level-docs:population.Population}]{\sphinxcrossref{\sphinxcode{\sphinxupquote{Population}}}}} model into file outf as a cPickle dump

\end{fulllineitems}



\section{{\hyperref[\detokenize{module-level-docs:module-radialmodels}]{\sphinxcrossref{\sphinxstyleliteralintitle{\sphinxupquote{radialmodels}}}}} \textendash{} Container for radial distn models}
\label{\detokenize{module-level-docs:module-radialmodels}}\label{\detokenize{module-level-docs:radialmodels-container-for-radial-distn-models}}\label{\detokenize{module-level-docs:radialmodels-module-docs}}\index{module@\spxentry{module}!radialmodels@\spxentry{radialmodels}}\index{radialmodels@\spxentry{radialmodels}!module@\spxentry{module}}\index{RadialModels (class in radialmodels)@\spxentry{RadialModels}\spxextra{class in radialmodels}}

\begin{fulllineitems}
\phantomsection\label{\detokenize{module-level-docs:radialmodels.RadialModels}}
\pysigstartsignatures
\pysigline
{\sphinxbfcode{\sphinxupquote{\DUrole{k}{class}\DUrole{w}{ }}}\sphinxcode{\sphinxupquote{radialmodels.}}\sphinxbfcode{\sphinxupquote{RadialModels}}}
\pysigstopsignatures
\end{fulllineitems}

\index{seed() (in module radialmodels)@\spxentry{seed()}\spxextra{in module radialmodels}}

\begin{fulllineitems}
\phantomsection\label{\detokenize{module-level-docs:radialmodels.seed}}
\pysigstartsignatures
\pysiglinewithargsret
{\sphinxcode{\sphinxupquote{radialmodels.}}\sphinxbfcode{\sphinxupquote{seed}}}
{}
{}
\pysigstopsignatures
\sphinxAtStartPar
Call the FORTRAN routine to make a seed

\end{fulllineitems}

\index{slabdist() (in module radialmodels)@\spxentry{slabdist()}\spxextra{in module radialmodels}}

\begin{fulllineitems}
\phantomsection\label{\detokenize{module-level-docs:radialmodels.slabdist}}
\pysigstartsignatures
\pysiglinewithargsret
{\sphinxcode{\sphinxupquote{radialmodels.}}\sphinxbfcode{\sphinxupquote{slabdist}}}
{}
{}
\pysigstopsignatures
\sphinxAtStartPar
Pick a point from a “slab” distribution around the Galactic plane

\end{fulllineitems}

\index{diskdist() (in module radialmodels)@\spxentry{diskdist()}\spxextra{in module radialmodels}}

\begin{fulllineitems}
\phantomsection\label{\detokenize{module-level-docs:radialmodels.diskdist}}
\pysigstartsignatures
\pysiglinewithargsret
{\sphinxcode{\sphinxupquote{radialmodels.}}\sphinxbfcode{\sphinxupquote{diskdist}}}
{}
{}
\pysigstopsignatures
\sphinxAtStartPar
Pick a point from a distribution purely along the Galactic plane

\end{fulllineitems}

\index{lfl06() (in module radialmodels)@\spxentry{lfl06()}\spxextra{in module radialmodels}}

\begin{fulllineitems}
\phantomsection\label{\detokenize{module-level-docs:radialmodels.lfl06}}
\pysigstartsignatures
\pysiglinewithargsret
{\sphinxcode{\sphinxupquote{radialmodels.}}\sphinxbfcode{\sphinxupquote{lfl06}}}
{}
{}
\pysigstopsignatures
\sphinxAtStartPar
Pick a point from the Lorimer et al (2006) Galactic distribution

\end{fulllineitems}

\index{ykr() (in module radialmodels)@\spxentry{ykr()}\spxextra{in module radialmodels}}

\begin{fulllineitems}
\phantomsection\label{\detokenize{module-level-docs:radialmodels.ykr}}
\pysigstartsignatures
\pysiglinewithargsret
{\sphinxcode{\sphinxupquote{radialmodels.}}\sphinxbfcode{\sphinxupquote{ykr}}}
{}
{}
\pysigstopsignatures
\sphinxAtStartPar
Pick a point from the Yusifov \& Kucuk Galactic distribution

\end{fulllineitems}



\section{{\hyperref[\detokenize{module-level-docs:module-galacticops}]{\sphinxcrossref{\sphinxstyleliteralintitle{\sphinxupquote{galacticops}}}}} \textendash{} Container for functions relating to the Galaxy}
\label{\detokenize{module-level-docs:module-galacticops}}\label{\detokenize{module-level-docs:galacticops-container-for-functions-relating-to-the-galaxy}}\label{\detokenize{module-level-docs:galacticops-module-docs}}\index{module@\spxentry{module}!galacticops@\spxentry{galacticops}}\index{galacticops@\spxentry{galacticops}!module@\spxentry{module}}\index{GalacticOps (class in galacticops)@\spxentry{GalacticOps}\spxextra{class in galacticops}}

\begin{fulllineitems}
\phantomsection\label{\detokenize{module-level-docs:galacticops.GalacticOps}}
\pysigstartsignatures
\pysigline
{\sphinxbfcode{\sphinxupquote{\DUrole{k}{class}\DUrole{w}{ }}}\sphinxcode{\sphinxupquote{galacticops.}}\sphinxbfcode{\sphinxupquote{GalacticOps}}}
\pysigstopsignatures
\end{fulllineitems}

\index{calc\_dtrue() (in module galacticops)@\spxentry{calc\_dtrue()}\spxextra{in module galacticops}}

\begin{fulllineitems}
\phantomsection\label{\detokenize{module-level-docs:galacticops.calc_dtrue}}
\pysigstartsignatures
\pysiglinewithargsret
{\sphinxcode{\sphinxupquote{galacticops.}}\sphinxbfcode{\sphinxupquote{calc\_dtrue}}}
{\sphinxparam{\DUrole{n}{(x}}\sphinxparamcomma \sphinxparam{\DUrole{n}{y}}\sphinxparamcomma \sphinxparam{\DUrole{n}{z)}}}
{}
\pysigstopsignatures
\sphinxAtStartPar
Calculate the distance from the Sun to Galactic coords (x, y, z) (NB. tuple)

\end{fulllineitems}

\index{calcXY() (in module galacticops)@\spxentry{calcXY()}\spxextra{in module galacticops}}

\begin{fulllineitems}
\phantomsection\label{\detokenize{module-level-docs:galacticops.calcXY}}
\pysigstartsignatures
\pysiglinewithargsret
{\sphinxcode{\sphinxupquote{galacticops.}}\sphinxbfcode{\sphinxupquote{calcXY}}}
{\sphinxparam{\DUrole{n}{r0}}}
{}
\pysigstopsignatures
\sphinxAtStartPar
Given a Galactic radius r0, choose an (x, y) position at random \(\theta\)

\end{fulllineitems}

\index{ne2025\_dist\_to\_dm() (in module galacticops)@\spxentry{ne2025\_dist\_to\_dm()}\spxextra{in module galacticops}}

\begin{fulllineitems}
\phantomsection\label{\detokenize{module-level-docs:galacticops.ne2025_dist_to_dm}}
\pysigstartsignatures
\pysiglinewithargsret
{\sphinxcode{\sphinxupquote{galacticops.}}\sphinxbfcode{\sphinxupquote{ne2025\_dist\_to\_dm}}}
{\sphinxparam{\DUrole{n}{dist}}\sphinxparamcomma \sphinxparam{\DUrole{n}{gl}}\sphinxparamcomma \sphinxparam{\DUrole{n}{gb}}}
{}
\pysigstopsignatures
\sphinxAtStartPar
Given a distance and Galactic coordinates, calculate DM according to NE2025

\end{fulllineitems}

\index{ne2001\_dist\_to\_dm() (in module galacticops)@\spxentry{ne2001\_dist\_to\_dm()}\spxextra{in module galacticops}}

\begin{fulllineitems}
\phantomsection\label{\detokenize{module-level-docs:galacticops.ne2001_dist_to_dm}}
\pysigstartsignatures
\pysiglinewithargsret
{\sphinxcode{\sphinxupquote{galacticops.}}\sphinxbfcode{\sphinxupquote{ne2001\_dist\_to\_dm}}}
{\sphinxparam{\DUrole{n}{dist}}\sphinxparamcomma \sphinxparam{\DUrole{n}{gl}}\sphinxparamcomma \sphinxparam{\DUrole{n}{gb}}}
{}
\pysigstopsignatures
\sphinxAtStartPar
Given a distance and Galactic coordinates, calculate DM according to NE2001

\end{fulllineitems}

\index{ymw16\_dist\_to\_dm() (in module galacticops)@\spxentry{ymw16\_dist\_to\_dm()}\spxextra{in module galacticops}}

\begin{fulllineitems}
\phantomsection\label{\detokenize{module-level-docs:galacticops.ymw16_dist_to_dm}}
\pysigstartsignatures
\pysiglinewithargsret
{\sphinxcode{\sphinxupquote{galacticops.}}\sphinxbfcode{\sphinxupquote{ymw16\_dist\_to\_dm}}}
{\sphinxparam{\DUrole{n}{dist}}\sphinxparamcomma \sphinxparam{\DUrole{n}{gl}}\sphinxparamcomma \sphinxparam{\DUrole{n}{gb}}}
{}
\pysigstopsignatures
\sphinxAtStartPar
Given a distance and Galactic coordinates, calculate DM according to the YMW16 model

\end{fulllineitems}

\index{lb\_to\_radec() (in module galacticops)@\spxentry{lb\_to\_radec()}\spxextra{in module galacticops}}

\begin{fulllineitems}
\phantomsection\label{\detokenize{module-level-docs:galacticops.lb_to_radec}}
\pysigstartsignatures
\pysiglinewithargsret
{\sphinxcode{\sphinxupquote{galacticops.}}\sphinxbfcode{\sphinxupquote{lb\_to\_radec}}}
{\sphinxparam{\DUrole{n}{gl}}\sphinxparamcomma \sphinxparam{\DUrole{n}{gb}}}
{}
\pysigstopsignatures
\sphinxAtStartPar
Convert Galactic coordinates to equatorial

\end{fulllineitems}

\index{ra\_dec\_to\_lb() (in module galacticops)@\spxentry{ra\_dec\_to\_lb()}\spxextra{in module galacticops}}

\begin{fulllineitems}
\phantomsection\label{\detokenize{module-level-docs:galacticops.ra_dec_to_lb}}
\pysigstartsignatures
\pysiglinewithargsret
{\sphinxcode{\sphinxupquote{galacticops.}}\sphinxbfcode{\sphinxupquote{ra\_dec\_to\_lb}}}
{\sphinxparam{\DUrole{n}{ra}}\sphinxparamcomma \sphinxparam{\DUrole{n}{dec}}}
{}
\pysigstopsignatures
\sphinxAtStartPar
Convert equatorial coordinates to Galactic

\end{fulllineitems}

\index{tsky() (in module galacticops)@\spxentry{tsky()}\spxextra{in module galacticops}}

\begin{fulllineitems}
\phantomsection\label{\detokenize{module-level-docs:galacticops.tsky}}
\pysigstartsignatures
\pysiglinewithargsret
{\sphinxcode{\sphinxupquote{galacticops.}}\sphinxbfcode{\sphinxupquote{tsky}}}
{\sphinxparam{\DUrole{n}{gl}}\sphinxparamcomma \sphinxparam{\DUrole{n}{gb}}\sphinxparamcomma \sphinxparam{\DUrole{n}{freq}}}
{}
\pysigstopsignatures
\sphinxAtStartPar
Calculate sky temperature at observing frequency freq and at Galactic coordinates gl, gb according to Haslam et al

\end{fulllineitems}

\index{xyz\_to\_lb() (in module galacticops)@\spxentry{xyz\_to\_lb()}\spxextra{in module galacticops}}

\begin{fulllineitems}
\phantomsection\label{\detokenize{module-level-docs:galacticops.xyz_to_lb}}
\pysigstartsignatures
\pysiglinewithargsret
{\sphinxcode{\sphinxupquote{galacticops.}}\sphinxbfcode{\sphinxupquote{xyz\_to\_lb}}}
{\sphinxparam{\DUrole{n}{(x}}\sphinxparamcomma \sphinxparam{\DUrole{n}{y}}\sphinxparamcomma \sphinxparam{\DUrole{n}{z)}}}
{}
\pysigstopsignatures
\sphinxAtStartPar
Convert the tuple (x, y, z) to Galactic sky coordinates.

\sphinxAtStartPar
Returns l, b in degrees

\end{fulllineitems}

\index{lb\_to\_xyz() (in module galacticops)@\spxentry{lb\_to\_xyz()}\spxextra{in module galacticops}}

\begin{fulllineitems}
\phantomsection\label{\detokenize{module-level-docs:galacticops.lb_to_xyz}}
\pysigstartsignatures
\pysiglinewithargsret
{\sphinxcode{\sphinxupquote{galacticops.}}\sphinxbfcode{\sphinxupquote{lb\_to\_xyz}}}
{\sphinxparam{\DUrole{n}{l}}\sphinxparamcomma \sphinxparam{\DUrole{n}{b}}\sphinxparamcomma \sphinxparam{\DUrole{n}{dist}}}
{}
\pysigstopsignatures
\sphinxAtStartPar
Convert Galactic sky coordinates at a distance dist to x,y,z coordinates.

\sphinxAtStartPar
Returns position as a tuple

\end{fulllineitems}

\index{scatter\_bhat() (in module galacticops)@\spxentry{scatter\_bhat()}\spxextra{in module galacticops}}

\begin{fulllineitems}
\phantomsection\label{\detokenize{module-level-docs:galacticops.scatter_bhat}}
\pysigstartsignatures
\pysiglinewithargsret
{\sphinxcode{\sphinxupquote{galacticops.}}\sphinxbfcode{\sphinxupquote{scatter\_bhat}}}
{\sphinxparam{\DUrole{n}{dm}}\sphinxparamcomma \sphinxparam{\DUrole{n}{scatterindex}}\sphinxparamcomma \sphinxparam{\DUrole{n}{freq\_mhz}}}
{}
\pysigstopsignatures
\sphinxAtStartPar
Calculate the scatter time according to Bhat et al at. Frequency in MHz, pulsar with dispersion measure dm, and using a scattering spectral index of scatterindex.

\sphinxAtStartPar
Calculated as

\sphinxAtStartPar
\(\tau = -6.46 + 0.154 \log_{10}({\rm{dm}}) + 1.07\log_{10}({\rm{dm}})^2 + {\rm{scatterindex}} \times \log_{10}(\frac{\rm{freq\_mhz}}{1000})\)

\sphinxAtStartPar
and typically \({\rm{scatterindex}} = -3.86\) (but there is an option to vary it!)

\end{fulllineitems}



\section{{\hyperref[\detokenize{module-level-docs:module-evolve}]{\sphinxcrossref{\sphinxstyleliteralintitle{\sphinxupquote{evolve}}}}} \textendash{} Evolve a pulsar population object}
\label{\detokenize{module-level-docs:evolve-evolve-a-pulsar-population-object}}\label{\detokenize{module-level-docs:evolve-module-docs}}\index{module@\spxentry{module}!evolve@\spxentry{evolve}}\index{evolve@\spxentry{evolve}!module@\spxentry{module}}\index{EvolveException (class in evolve)@\spxentry{EvolveException}\spxextra{class in evolve}}\phantomsection\label{\detokenize{module-level-docs:module-evolve}}

\begin{fulllineitems}
\phantomsection\label{\detokenize{module-level-docs:evolve.EvolveException}}
\pysigstartsignatures
\pysigline
{\sphinxbfcode{\sphinxupquote{\DUrole{k}{class}\DUrole{w}{ }}}\sphinxcode{\sphinxupquote{evolve.}}\sphinxbfcode{\sphinxupquote{EvolveException}}}
\pysigstopsignatures
\end{fulllineitems}

\index{generate() (in module evolve)@\spxentry{generate()}\spxextra{in module evolve}}

\begin{fulllineitems}
\phantomsection\label{\detokenize{module-level-docs:evolve.generate}}
\pysigstartsignatures
\pysiglinewithargsret
{\sphinxcode{\sphinxupquote{evolve.}}\sphinxbfcode{\sphinxupquote{generate}}}
{\sphinxparam{\DUrole{n}{ngen}}\sphinxparamcomma \sphinxparam{\DUrole{n}{surveyList}\DUrole{o}{=}\DUrole{default_value}{None}}\sphinxparamcomma \sphinxparam{\DUrole{n}{age\_max}\DUrole{o}{=}\DUrole{default_value}{1.0E9}}\sphinxparamcomma \sphinxparam{\DUrole{n}{pDistPars}\DUrole{o}{=}\DUrole{default_value}{{[}.3, .15{]}}}\sphinxparamcomma \sphinxparam{\DUrole{n}{bFieldPars}\DUrole{o}{=}\DUrole{default_value}{{[}12.65, 0.55{]}}}\sphinxparamcomma \sphinxparam{\DUrole{n}{birthVPars}\DUrole{o}{=}\DUrole{default_value}{{[}0.0, 180.{]}}}\sphinxparamcomma \sphinxparam{\DUrole{n}{siDistPars}\DUrole{o}{=}\DUrole{default_value}{{[}\sphinxhyphen{}1.6, 0.35{]}}}\sphinxparamcomma \sphinxparam{\DUrole{n}{alignModel}\DUrole{o}{=}\DUrole{default_value}{\textquotesingle{}orthogonal\textquotesingle{}}}\sphinxparamcomma \sphinxparam{\DUrole{n}{lumDistType}\DUrole{o}{=}\DUrole{default_value}{\textquotesingle{}fk06\textquotesingle{}}}\sphinxparamcomma \sphinxparam{\DUrole{n}{lumDistPars}\DUrole{o}{=}\DUrole{default_value}{{[}\sphinxhyphen{}1.5, 0.5{]}}}\sphinxparamcomma \sphinxparam{\DUrole{n}{alignTime}\DUrole{o}{=}\DUrole{default_value}{None}}\sphinxparamcomma \sphinxparam{\DUrole{n}{spinModel}\DUrole{o}{=}\DUrole{default_value}{\textquotesingle{}fk06\textquotesingle{}}}\sphinxparamcomma \sphinxparam{\DUrole{n}{beamModel}\DUrole{o}{=}\DUrole{default_value}{\textquotesingle{}tm98\textquotesingle{}}}\sphinxparamcomma \sphinxparam{\DUrole{n}{birthVModel}\DUrole{o}{=}\DUrole{default_value}{\textquotesingle{}gaussian\textquotesingle{}}}\sphinxparamcomma \sphinxparam{\DUrole{n}{electronModel}\DUrole{o}{=}\DUrole{default_value}{\textquotesingle{}ne2025\textquotesingle{}}}\sphinxparamcomma \sphinxparam{\DUrole{n}{braking\_index}\DUrole{o}{=}\DUrole{default_value}{0}}\sphinxparamcomma \sphinxparam{\DUrole{n}{zscale}\DUrole{o}{=}\DUrole{default_value}{0.05}}\sphinxparamcomma \sphinxparam{\DUrole{n}{duty}\DUrole{o}{=}\DUrole{default_value}{5.}}\sphinxparamcomma \sphinxparam{\DUrole{n}{scindex}\DUrole{o}{=}\DUrole{default_value}{\sphinxhyphen{}3.86}}\sphinxparamcomma \sphinxparam{\DUrole{n}{widthModel}\DUrole{o}{=}\DUrole{default_value}{None}}\sphinxparamcomma \sphinxparam{\DUrole{n}{nodeathline}\DUrole{o}{=}\DUrole{default_value}{False}}\sphinxparamcomma \sphinxparam{\DUrole{n}{efficiencycut}\DUrole{o}{=}\DUrole{default_value}{None}}\sphinxparamcomma \sphinxparam{\DUrole{n}{nostdout}\DUrole{o}{=}\DUrole{default_value}{False}}\sphinxparamcomma \sphinxparam{\DUrole{n}{nospiralarms}\DUrole{o}{=}\DUrole{default_value}{False}}\sphinxparamcomma \sphinxparam{\DUrole{n}{keepdead}\DUrole{o}{=}\DUrole{default_value}{False}}\sphinxparamcomma \sphinxparam{\DUrole{n}{ascfile}\DUrole{o}{=}\DUrole{default_value}{None}}}
{}
\pysigstopsignatures
\sphinxAtStartPar
Generate an evolved pulsar population model.

\end{fulllineitems}

\index{main() (in module evolve)@\spxentry{main()}\spxextra{in module evolve}}

\begin{fulllineitems}
\phantomsection\label{\detokenize{module-level-docs:evolve.main}}
\pysigstartsignatures
\pysiglinewithargsret
{\sphinxcode{\sphinxupquote{evolve.}}\sphinxbfcode{\sphinxupquote{main}}}
{}
{}
\pysigstopsignatures
\sphinxAtStartPar
Entry point for the evolve console script.

\end{fulllineitems}



\chapter{Indices and tables}
\label{\detokenize{index:indices-and-tables}}\begin{itemize}
\item {} 
\sphinxAtStartPar
\DUrole{xref}{\DUrole{std}{\DUrole{std-ref}{genindex}}}

\item {} 
\sphinxAtStartPar
\DUrole{xref}{\DUrole{std}{\DUrole{std-ref}{modindex}}}

\item {} 
\sphinxAtStartPar
\DUrole{xref}{\DUrole{std}{\DUrole{std-ref}{search}}}

\end{itemize}


\renewcommand{\indexname}{Python Module Index}
\begin{sphinxtheindex}
\let\bigletter\sphinxstyleindexlettergroup
\bigletter{e}
\item\relax\sphinxstyleindexentry{evolve}\sphinxstyleindexpageref{module-level-docs:\detokenize{module-evolve}}
\indexspace
\bigletter{g}
\item\relax\sphinxstyleindexentry{galacticops}\sphinxstyleindexpageref{module-level-docs:\detokenize{module-galacticops}}
\indexspace
\bigletter{p}
\item\relax\sphinxstyleindexentry{populate}\sphinxstyleindexpageref{module-level-docs:\detokenize{module-populate}}
\item\relax\sphinxstyleindexentry{population}\sphinxstyleindexpageref{module-level-docs:\detokenize{module-population}}
\item\relax\sphinxstyleindexentry{pulsar}\sphinxstyleindexpageref{module-level-docs:\detokenize{module-pulsar}}
\indexspace
\bigletter{r}
\item\relax\sphinxstyleindexentry{radialmodels}\sphinxstyleindexpageref{module-level-docs:\detokenize{module-radialmodels}}
\indexspace
\bigletter{s}
\item\relax\sphinxstyleindexentry{survey}\sphinxstyleindexpageref{module-level-docs:\detokenize{module-survey}}
\end{sphinxtheindex}

\renewcommand{\indexname}{Index}
\printindex
\end{document}